% ============ 第 IV 部分: \index{Peter-Weyl 定理}Peter-Weyl 定理 (The Peter-Weyl Theorem) ============
% 原书 IV.1 - IV.9 (原书 pp.195-229)
% 编号: IV.1.1 (节级) / 习题 IV.9.1 (节级)

\section{\index{紧算子}紧算子}

\begin{definition}[紧算子 (compact operator)]
\label{def:IV.1.1}
设 $V$ 和 $W$ 是巴拿赫空间, $B$ 是 $V$ 中的闭单位球. 若 $T(B)$ 紧, 则称连续线性算子 $T \colon V \to W$ 紧 (compact).
\end{definition}

\begin{example}[紧算子的例子]
\label{exa:IV.1.2}
\begin{enumerate}
\item[(1)] 若 $\Im(T)$ 有限维 (即 $T$ 有有限秩 (finite rank)), 则 $T$ 紧.
\item[(2)] 若 $T$ 是有限秩算子的极限, 则 $T$ 紧; 更一般地, 紧算子的任何极限都紧 (见习题 \ref{ex:I.5.5.9}).
  反之, 若 $W$ 是希尔伯特空间, 则每个紧算子 $T \colon V \to W$ 都是有限秩算子的极限.\footnote{一般情形这不成立, 见 Enflo 的文章 {[}10{]} 中的反例.}
\item[(3)] $V$ 的恒等算子紧当且仅当 $V$ 有限维. (这是 Riesz 引理的推论, 见定理 \ref{thm:B.4.2}.)
\end{enumerate}
\end{example}

在本课程中, 我们只需使用希尔伯特空间的自伴紧自同态. 对它们有谱定理的一个更简单的版本.

\begin{theorem}[\index{自伴紧算子的谱定理}自伴紧算子的谱定理]
\label{thm:IV.1.3}
设 $V$ 是 $\CC$ 上的希尔伯特空间, $T \colon V \to V$ 是 $V$ 的连续自同态. 假设 $T$ 紧且自伴, 对每个 $\lambda \in \CC$ 记 $V_{\lambda} = \Ker(T - \lambda\mathrm{id}_V)$.
则:
\begin{enumerate}
\item[(i)] 若 $V_{\lambda} \neq 0$, 则 $\lambda \in \R$.
\item[(ii)] 若 $\lambda, \mu \in \CC$ 且 $\lambda \neq \mu$, 则 $V_{\mu} \subset V_{\lambda}^{\perp}$.
\item[(iii)] 若 $\lambda \in \CC - \{0\}$, 则 $\dim_{\CC}V_{\lambda} < +\infty$.
\item[(iv)] $\{\lambda \in \CC \mid V_{\lambda} \neq 0\}$ 有限或可数, 且其唯一可能的极限点是 $0$.
\item[(v)] $\bigoplus_{\lambda \in \CC}V_{\lambda}$ 在 $V$ 中稠密.
\end{enumerate}
\end{theorem}

\begin{proof}
我们证明 (i). 设 $\lambda \in \CC$ 使得 $V_{\lambda} \neq 0$, 选取非零的 $v \in V_{\lambda}$. 则
\[
\lambda\|v\|^2 = \langle \lambda v, v \rangle = \langle T(v), v \rangle = \langle v, T^*(v) \rangle = \langle v, T(v) \rangle = \bar{\lambda}\|v\|^2.
\]
由于 $\|v\| \neq 0$, 这蕴含 $\lambda \in \R$.
我们证明 (ii). 设 $\lambda, \mu \in \CC$ 满足 $\lambda \neq \mu$, $v \in V_{\lambda}$, $w \in V_{\mu}$. 要证明 $\langle v, w \rangle = 0$. 由 (i), 只需处理 $\lambda, \mu \in \R$ 的情形 (否则 $v = w = 0$). 此时
\[
\lambda\langle v, w \rangle = \langle T(v), w \rangle = \langle v, T(w) \rangle = \mu\langle v, w \rangle,
\]
所以 $\langle v, w \rangle = 0$.
设 $r > 0$. 设 $W = \bigoplus_{|\lambda| \geq r}V_{\lambda}$. 要证明 $\dim W < +\infty$, 这将蕴含 (iii) 和 (iv). 选取 $W$ 的由 $T$ 的特征向量组成的 Hilbert 基 $(e_i)_{i \in I}$, 即对每个 $i \in I$, $T(e_i) = \lambda_ie_i$, 其中 $|\lambda_i| \geq r$. 若 $I$ 无穷, 则族 $(T(e_i))_{i \in I}$ 不可能有收敛的 (非稳定的) 子序列. 事实上, 若存在单射映射 $\N \to I$, $n \longmapsto i_n$, 使得 $(T(e_{i_n}))_{n \geq 0}$ 收敛到 $V$ 的某个向量 $v$, 则 $\lambda_{i_n}e_{i_n} \to v$, 所以 $v$ 在 $\operatorname{Span}(e_{i_n} \mid n \geq 0)$ 的闭包中. 但另一方面, 对每个 $n \geq 0$,
\[
\langle v, e_{i_n} \rangle = \lim_{m \to +\infty}\langle \lambda_{i_m}e_{i_m}, e_{i_n} \rangle = 0,
\]
所以 $v \in \operatorname{Span}(e_{i_n} \mid n \geq 0)^{\perp}$. 这迫使 $v = 0$. 但 $\|v\| = \lim_{n \to +\infty}\|\lambda_{i_n}e_{i_n}\| \geq r > 0$, 矛盾. 由于 $T$ 紧, 这证明 $I$ 不能无穷, 即 $\dim(W) < +\infty$.
让我们证明 (v). 设 $W' = \bigoplus_{\lambda \in \CC}V_{\lambda}$, $W = W'^{\perp}$. 要证明 $W = 0$. 所以假设 $W \neq 0$. 由于 $T$ 自伴且 $W'$ 显然被 $T$ 稳定, 有 $T(W) \subset W$. (若 $v \in W$, 则对每个 $w \in W'$, $\langle T(v), w \rangle = \langle v, T(w) \rangle = 0$.) 由 $W$ 的定义, $\Ker(T|_W) = \{0\}$, 所以 $\|T|_W\|_{\mathrm{op}} > 0$. 设 $B = \{x \in W \mid \|x\| = 1\}$. 由于由下面的引理, $\|T|_W\|_{\mathrm{op}} = \sup_{x \in B}|\langle T(x), x \rangle|$, 存在 $B$ 的元素序列 $(x_n)_{n \geq 0}$ 使得当 $n \to +\infty$ 时 $\langle T(x_n), x_n \rangle \to \lambda$, 其中 $\lambda = \pm\|T|_W\|_{\mathrm{op}}$. 则
\[
0 \leq \|T(x_n) - \lambda x_n\|^2 = \|T(x_n)\|^2 + \lambda^2\|x_n\|^2 - 2\lambda\langle T(x_n), x_n \rangle \leq 2\lambda^2 - 2\lambda\langle T(x_n), x_n \rangle
\]
在 $n \to +\infty$ 时收敛到 $0$, 所以 $T(x_n) - \lambda x_n$ 本身收敛到 $0$. 由于 $T$ 紧, 可以假设序列 $(T(x_n))_{n \geq 0}$ 在 $W$ 中有极限, 设为 $w$. 则 $T(w) - \lambda w = 0$. 由 $W$ 的定义, 必有 $w = 0$. 但此时 $T(x_n) \to 0$, 所以 $\langle T(x_n), x_n \rangle \to 0$, 从而 $\lambda = 0 = \|T|_W\|_{\mathrm{op}}$, 矛盾.
\end{proof}

\begin{lemma}[自伴算子的范数]
\label{lem:IV.1.4}
设 $V$ 是希尔伯特空间, $T \in \End(V)$ 自伴. 则
\[
\|T\|_{\mathrm{op}} = \sup_{x \in V, \|x\| = 1}|\langle T(x), x \rangle|.
\]
\end{lemma}

\begin{proof}
设 $c = \sup_{x \in V, \|x\| = 1}|\langle T(x), x \rangle|$. 由 $\|T\|_{\mathrm{op}}$ 的定义有 $c \leq \|T\|_{\mathrm{op}}$. 由于 $\|T\|_{\mathrm{op}}$ 等于在范数 $1$ 向量上的 $\sup_{x,y \in V}|\langle T(x), y \rangle|$, 为证另一个不等式只需证明对所有满足 $\|x\| = \|y\| = 1$ 的 $x, y \in V$ 有 $|\langle T(x), y \rangle| \leq c$. 设 $x, y \in V$. 用 $\CC$ 的范数 $1$ 的元素乘 $y$ (这不改变 $\|y\|$) 后, 可以假设 $\langle T(x), y \rangle \in \R$. 则
\[
\langle T(x), y \rangle = \frac{1}{4}(\langle T(x + y), x + y \rangle - \langle T(x - y), x - y \rangle).
\]
所以, 利用平行四边形恒等式,
\begin{multline*}
|\langle T(x), y \rangle| \leq \frac{c}{4}(\|x + y\|^2 + \|x - y\|^2) \\
= \frac{c}{2}(\|x\|^2 + \|y\|^2).
\end{multline*}
这证明所要的结果.
\end{proof}

这里有一些对更一般情形的紧算子成立的结果 (见 {[}20{]} 4.16-4.25).

\begin{theorem}[紧算子谱的性质]
\label{thm:IV.1.5}
设 $V$ 是巴拿赫空间, $T \in \End(V)$ 是紧自同态. 记 $\sigma(T)$ 为 $T$ 在 $\End(V)$ 中的谱, 即
\[
\sigma(T) = \{\lambda \in \CC \mid \lambda\mathrm{id}_V - T \notin \End(V)^{\times}\}.
\]
则:
\begin{enumerate}
\item[(i)] 对每个 $\lambda \neq 0$, $T - \lambda\mathrm{id}_V$ 的像闭.
\item[(ii)] 对每个 $\lambda \in \sigma(T) - \{0\}$, 有 $\Im(T - \lambda\mathrm{id}_V) \neq V$ 且 $\dim(\Ker(T - \lambda\mathrm{id}_V)) = \dim(V/\Im(T - \lambda\mathrm{id}_V))$.\footnote{注意这推广了秩-零度定理.} 特别地, $\Ker(T - \lambda\mathrm{id}_V) \neq \{0\}$.
\item[(iii)] 对每个 $\lambda \neq 0$, 递增序列 $(\Ker((T - \lambda\mathrm{id}_V)^n))_{n \geq 1}$ 稳定, 且其极限有限维.
\item[(iv)] 若 $\dim_{\CC}V = +\infty$, 则 $0 \in \sigma(T)$.
\item[(v)] $\CC - \{0\}$ 的子集 $\sigma(T) - \{0\}$ 离散. 特别地, 对每个 $r > 0$, 只有有限多个 $\lambda \in \sigma(T)$ 满足 $|\lambda| \geq r$.
\end{enumerate}
特别地, 若 $V$ 是希尔伯特空间且 $T$ 自伴, 则定理 \ref{thm:IV.1.3} 的 (v) 变成
\[
V = \bigoplus_{\lambda \in \sigma(T)}\Ker(T - \lambda\mathrm{id}_V).
\]
\end{theorem}

\section{紧群的酉表示的半单性}
本节的目标是证明下面的定理. (与命题 \ref{prop:I.3.3.3} 比较.)

\begin{theorem}[紧群酉表示的半单性]
\label{thm:IV.2.1}
设 $G$ 是紧群, $V$ 是 $G$ 的酉表示. 则存在 $V$ 的两两正交的子表示族 $(W_i)_{i \in I}$ 使得每个 $W_i$ 不可约, 且
\[
V = \bigoplus_{i \in I}W_i.
\]
\end{theorem}

我们已经在习题 \ref{ex:I.5.5.9} 中见过关键的构造. 让我们把它总结成一个命题.

\begin{proposition}[等变自伴紧算子]
\label{prop:IV.2.2}
(见习题 \ref{ex:I.5.5.9}.) 设 $G$ 是紧群, $dx$ 是 $G$ 上的规范化 Haar 测度, $(\pi, V)$ 是 $G$ 的酉表示. 若 $u \in V$, 则公式
\[
T(v) = \int_G \langle v, \pi(x)(u) \rangle\pi(x)(u)\,\dd x
\]
定义 $V$ 的连续 $G$-等变自伴紧自同态, 且 $T = 0$ 当且仅当 $u = 0$.
\end{proposition}

事实上, 我们甚至知道 $T$ 正定, 即对每个 $v \in V$ 有 $\langle T(v), v \rangle \geq 0$.

\begin{corollary}[不可约子表示的存在性]
\label{cor:IV.2.3}
设 $V$ 是紧群 $G$ 的非零酉表示. 则 $V$ 包含 $G$ 的不可约表示.
\end{corollary}

\begin{proof}
若 $V$ 有限维, 则 $V$ 的任何极小维数的非零 $G$-不变子空间必不可约.
一般情形下, 选取 $u \in V - \{0\}$, 设 $T \in \End(V)$ 是命题中构造的 $V$ 的自同态. 由自伴紧算子的谱定理 (定理 \ref{thm:IV.1.3}),
\[
V = \bigoplus_{\lambda \in \CC}\Ker(T - \lambda\mathrm{id}_V).
\]
由于 $T \neq 0$, $V$ 的闭子空间 $\Ker(T)$ 不等于 $V$. 由上面的等式, 存在 $\lambda \in \CC - \{0\}$ 使得 $W := \Ker(T - \lambda\mathrm{id}_V) \neq 0$. 则 $W$ 是 $V$ 的非零闭子空间, 且它 $G$-不变 (因为 $T$ 是 $G$-等变的) 且由定义被 $T$ 稳定. 又, 由定理 \ref{thm:IV.1.3}(iii), 空间 $W$ 有限维. 所以由证明的开头, $W$ 有不可约子表示, 我们完成了.
\end{proof}

\begin{proof}[定理的证明]
由 Zorn 引理, 可以找到 $V$ 的两两正交的不可约子表示的极大族 $(W_i)_{i \in I}$. 假设 $W_i$ 的直和不在 $V$ 中稠密, 则 $W := \bigl(\bigoplus_{i \in I}W_i\bigr)^{\perp}$ 是 $V$ 的非零闭不变子空间 (见引理 \ref{lem:I.3.2.6}). 由上面的推论, 表示 $W$ 有不可约子表示, 这与族 $(W_i)_{i \in I}$ 的极大性矛盾. 所以 $V = \bigoplus_{i \in I}W_i$.
\end{proof}

我们用关于酉表示的两种等价概念的一个注记结束本节. 回忆拓扑群 $G$ 的两个连续表示 $V_1$ 和 $V_2$ 称为等价 (equivalent) (或同构 (isomorphic)), 如果存在有连续逆的连续 $G$-等变同构 $V_1 \to V_2$.

\begin{definition}[\index{酉等价}酉等价 (unitarily equivalent)]
\label{def:IV.2.4}
拓扑群 $G$ 的两个酉表示 $V_1$ 和 $V_2$ 称为酉等价, 如果存在是等距的 $G$-等变同构 $V_1 \to V_2$.
\end{definition}

两个酉等价的表示显然等价.

\begin{example}[左右正则表示的酉等价]
\label{exa:IV.2.5}
设 $G$ 是局部紧群, $\mu$ 是 $G$ 上的左 Haar 测度, $\nu$ 是由 $\nu(E) = \mu(E^{-1})$ 定义的右 Haar 测度.
则 $G$ 的左、右正则表示酉等价, 同构把 $f \in L^2(G, \mu)$ 送到 $L^2(G, \nu)$ 的元素 $x \longmapsto \Delta(x)^{-1/2}\overline{f(x^{-1})}$. (见命题 \ref{prop:I.2.12}.)
\end{example}

\begin{proposition}[不可约表示的等价与酉等价]
\label{prop:IV.2.6}
假设 $V_1$ 和 $V_2$ 是 $G$ 的不可约酉表示. 则它们等价当且仅当它们酉等价.
\end{proposition}

\begin{proof}
假设 $V_1$ 和 $V_2$ 等价, 设 $U \colon V_1 \to V_2$ 是 $G$-等变同构. 记 $\langle \cdot, \cdot \rangle_1$ 和 $\langle \cdot, \cdot \rangle_2$ 为 $V_1$ 和 $V_2$ 的内积. 设 $B \colon V_1 \times V_1 \to \CC$, $(v, w) \longmapsto \langle U(v), U(w) \rangle_2$. 这是 $V_1$ 上的 Hermitian 半双线性形式, 且因 $U$ 有界而有界. 由下面的引理, 存在自伴自同态 $T \in \End(V_1)$ 使得对所有 $v, w \in V_1$ 有 $B(v, w) = \langle T(v), w \rangle_1$. 让我们证明 $T$ 是 $G$-等变的. 设 $v \in V_1$, $x \in G$. 对每个 $w \in V_1$,
\begin{align*}
\langle T(\pi_1(x)(v)), w \rangle_1 &= B(\pi_1(x)(v), w) \\
&= \langle U(\pi_1(x)(v)), U(w) \rangle_2 \\
&= \langle \pi_2(x)(U(v)), U(w) \rangle_2 \\
&= \langle U(v), \pi_2(x^{-1})U(w) \rangle_2 \\
&= \langle U(v), U(\pi_1(x)^{-1}(w)) \rangle_2 \\
&= B(v, \pi_1(x^{-1})(w)) \\
&= \langle T(v), \pi_1(x)^{-1}(w) \rangle_1 \\
&= \langle \pi_1(x)(T(v)), w \rangle_1,
\end{align*}
所以 $T(\pi_1(x)(v)) = \pi_1(x)(T(v))$. 由于 $V_1$ 不可约, Schur 引理 (定理 \ref{thm:I.3.4.1}) 蕴含对某个 $\lambda \in \CC$ 有 $T = \lambda\mathrm{id}_{V_1}$. 由于对每个非零 $v \in V_1$ 有 $\langle T(v), v \rangle_1 = \langle U(v), U(v) \rangle_2 > 0$, 必有 $\lambda \in \R_{>0}$. 则 $\lambda^{-1/2}U$ 是等距, 所以 $V_1$ 和 $V_2$ 酉等价.
\end{proof}

\begin{lemma}[有界半双线性形式]
\label{lem:IV.2.7}
设 $V$ 是希尔伯特空间, $B \colon V \times V \to \CC$ 是有界半双线性形式 (即 $B$ 对第一个变量 $\CC$-线性、对第二个变量 $\CC$-反线性; 有界性条件指 $\sup_{v,w \in V, \|v\| = \|w\| = 1}|B(v, w)| < +\infty$).
则存在唯一的 $T \in \End(V)$ 使得对所有 $v, w \in V$,
\[
B(v, w) = \langle T(v), w \rangle.
\]
而且, $T$ 自伴当且仅当 $B$ Hermitian (即对所有 $v, w \in V$ 有 $B(w, v) = \overline{B(v, w)}$).
\end{lemma}

\begin{proof}
$T$ 的唯一性显然 (它由 $V^{\perp} = \{0\}$ 这一事实得出.)
若 $v \in V$, 则映射 $V \to \CC$, $w \longmapsto B(v, w)$ 是 $V$ 上的连续线性泛函, 所以存在唯一的 $T(v) \in V$ 使得对每个 $w \in V$ 有 $B(v, w) = \langle T(v), w \rangle$. $T$ 的线性由 $B$ 对第一个变量线性得出. 而且, 对每个 $v \in V$,
\[
\|T(v)\| = \sup_{w \in V, \|w\| = 1}|\langle T(v), w \rangle| = \sup_{w \in V, \|w\| = 1}|B(v, w)| \leq C\|v\|,
\]
其中
\[
C = \sup_{x,y \in V, \|x\| = \|y\| = 1}|B(x, y)|.
\]
所以 $T$ 有界.
最后, $T$ 自伴当且仅当对所有 $v, w \in V$,
\[
B(v, w) = \langle T(v), w \rangle = \langle v, T(w) \rangle = \overline{B(w, v)}.
\]
这证明最后一个结论.
\end{proof}

\begin{definition}[\index{酉对偶}酉对偶 (unitary dual)]
\label{def:IV.2.8}
记 $\widehat{G}$ 为 $G$ 的不可约酉表示的等价 (或酉等价) 类的集合, 称之为 $G$ 的酉对偶.
若 $(\pi, V) \in \widehat{G}$, 记 $\dim(\pi)$ 和 $\End(\pi)$ 为 $\dim(V)$ 和 $\End(\pi)$.
\end{definition}

注意此记号与习题 \ref{ex:I.5.4.1} 中对交换群使用的记号一致.

\section{\index{矩阵系数}矩阵系数}

\begin{definition}[矩阵系数 (matrix coefficient)]
\label{def:IV.3.1}
设 $(\pi, V)$ 是拓扑群 $G$ 的酉表示. $(\pi, V)$ 的矩阵系数是形如 $x \longmapsto \langle \pi(x)(u), v \rangle$ ($u, v \in V$) 的函数 $G \to \CC$.
\end{definition}

注意矩阵系数是连续函数. 记 $E_{\pi}$ 或 $E_V$ 为 $\mathcal{C}(G)$ 中由 $\pi$ 的矩阵系数张成的子空间.
我们先证明对任意群 $G$ 成立的一些一般结果.

\begin{proposition}[矩阵系数的基本性质]
\label{prop:IV.3.2}
设 $(\pi, V)$ 是 $G$ 的酉表示.
\begin{enumerate}
\item[(i)] $\mathcal{C}(G)$ 的子空间 $E_{\pi}$ 只依赖 $\pi$ 的酉等价类, 且对每个 $x \in G$ 被算子 $L_x$ 和 $R_x$ 不变.
\item[(ii)] 若 $V$ 有限维, 则 $E_{\pi}$ 有限维且 $\dim(E_{\pi}) \leq (\dim V)^2$.
\item[(iii)] 若 $V = V_1 \oplus \cdots \oplus V_n$ 且 $V_i$ 两两正交且 $G$-不变, 则 $E_{\pi} = \sum_{i=1}^n E_{V_i}$.
\item[(iv)] 有 $E_{\pi^*} = \overline{E_{\pi}}$.
\end{enumerate}
特别地, 我们得到 $G \times G$ 在 $E_{\pi}$ 上的作用, 由 $(x, y) \in G \times G$ 按 $L_x \circ R_y = R_y \circ L_x$ 作用.
\end{proposition}

\begin{proof}
(i). 第一个结论显然. 为证明第二个结论, 设 $v, w \in V$, $x \in G$. 则对每个 $y \in G$,
\[
\langle \pi(x^{-1}y)(v), w \rangle = \langle \pi(y)(v), \pi(x)(w) \rangle
\]
且
\[
\langle \pi(yx)(v), w \rangle = \langle \pi(y)(\pi(x)(v)), w \rangle,
\]
所以函数 $y \longmapsto \langle \pi(x^{-1}y)(v), w \rangle$ 和 $y \longmapsto \langle \pi(yx)(v), w \rangle$ 也是 $\pi$ 的矩阵系数.
(ii). 设 $(e_1, \ldots, e_n)$ 是 $V$ 的基. 对 $i, j \in \{1, \ldots, n\}$, 记 $\phi_{ij}$ 为函数 $G \to \CC$, $x \longmapsto \langle \pi(x)(e_i), e_j \rangle$. 若 $v, w \in V$, 可以把 $v$ 写成 $v = \sum_{i=1}^n a_ie_i$, $w = \sum_{j=1}^n b_je_j$, 然后对每个 $x \in G$,
\[
\langle \pi(x)(v), w \rangle = \sum_{i,j=1}^n a_i\bar{b}_j\phi_{ij}(x).
\]
所以族 $(\phi_{ij})_{1 \leq i,j \leq n}$ 张成 $E_{\pi}$.
(iii). 对每个 $i \in \{1, \ldots, n\}$, 显然有 $E_{V_i} \subset E_{\pi}$. 所以 $\sum_{i=1}^n E_{V_i} \subset E_{\pi}$. 反之, 设 $v, w \in V$, 把 $v$ 写成 $v = \sum_{i=1}^n v_i$, $w = \sum_{i=1}^n w_i$, 其中 $v_i, w_i \in V_i$. 则对每个 $x \in G$,
\[
\langle \pi(x)(v), w \rangle = \sum_{i,j=1}^n\langle \pi(x)(v_i), w_j \rangle = \sum_{i=1}^n\langle \pi(x)(v_i), w_i \rangle.
\]
所以函数 $x \longmapsto \langle \pi(x)(v), w \rangle$ 在 $\sum_{i=1}^n E_{V_i}$ 中.
\end{proof}

\begin{definition}[$\Hom(V, V')$ 上的 $G \times G$ 作用]
\label{def:IV.3.3}
设 $(\pi, V)$ 和 $(\pi', V')$ 是 $V$ 的连续表示. 定义 $G \times G$ 在 $\Hom(V, V')$ 上的作用 $\rho$ 为
\[
\rho(x, y)(T) = \pi'(y) \circ T \circ \pi(x)^{-1},
\]
其中 $T \in \Hom(V, V')$, $x, y \in G$.
\end{definition}

\begin{proposition}[交错算子与不动点]
\label{prop:IV.3.4}
有
\[
\Hom_G(V, V') = \{T \in \Hom(V, V') \mid \forall x \in G,\, \rho(x, x)(T) = T\}.
\]
而且, 若映射 $G \to \End(V)$, $x \longmapsto \pi(x)$ 和 $G \to \End(V')$, $x \longmapsto \pi'(x)$ 连续 (例如若 $V$ 和 $V'$ 有限维, 见命题 \ref{prop:I.3.5.1}), 则上面定义的作用是 $G \times G$ 在 $\Hom(V, V')$ 上的连续表示.
\end{proposition}

\begin{proof}
第一个结论显然. 第二个结论由 $\Hom$ 空间上复合的连续性以及 $G$ 上取逆的连续性得出.
\end{proof}

特别地, 我们得到 $G \times G$ 在 $\End(V)$ 和 $V^* := \Hom(V, \CC)$ 上的作用 (用 $G$ 在 $\CC$ 上的平凡作用); 第二个限制到第一个因子上给出 $G$ 在 $V^*$ 上的作用 (若 $x \in G$, $\Lambda \in V^*$, 则 $(x, \Lambda)$ 被送到 $\Lambda \circ \pi(x)^{-1}$). 这将是这些空间上的默认作用.

\begin{definition}[张量积上的 $G \times G$ 作用]
\label{def:IV.3.5}
设 $(\pi, V)$ 和 $(\pi', V')$ 是 $V$ 的连续表示. 定义 $G \times G$ 在代数张量积 $V \otimes_{\CC}V'$ 上的作用 $\rho$ 为
\[
\rho(x, y)(v \otimes w) = \pi(x)(v) \otimes \pi'(y)(w),
\]
其中 $x, y \in G$, $v \in V$, $w \in V'$.
\end{definition}

这个作用良定义, 因为对所有 $x, y \in G$, 映射 $V \times V' \to V \otimes_{\CC}V'$, $(v, w) \longmapsto \pi(x)(v) \otimes \pi'(y)(w)$ 双线性, 所以它诱导映射 $\rho(x, y) \colon V \otimes_{\CC}V' \to V \otimes_{\CC}V'$. 若 $V$ 和 $V'$ 有限维, 则由命题 \ref{prop:I.3.5.1}, $G \times G$ 在 $V \otimes_{\CC}V'$ 上的作用连续.
注意, 若把 $G \times G$ 在 $V \otimes_{\CC}W$ 上的作用限制到第一个 (相应地第二个) 因子, 得到等价于 $V^{\oplus \dim(W)}$ (相应地 $W^{\oplus \dim(V)}$) 的表示.

\begin{proposition}[$V^* \otimes W$ 与 $\Hom(V, W)$]
\label{prop:IV.3.6}
设 $V, W$ 是 $G$ 的连续表示.
\begin{enumerate}
\item[(i)] 映射 $V^* \otimes_{\CC}W \to \Hom(V, W)$, 把 $\Lambda \otimes w$ ($\Lambda \in V^*$, $w \in W$) 送到线性算子 $V \to W$, $v \longmapsto \Lambda(v)w$ 良定义且 $G \times G$-等变. 若 $V$ 和 $W$ 有限维, 它是连续表示的等价.
\item[(ii)] 映射 $V^* \otimes_{\CC}V \to \mathcal{C}(G)$, 把 $\Lambda \otimes v$ ($\Lambda \in V^*$, $v \in V$) 送到函数 $G \to \CC$, $x \longmapsto \Lambda(\pi(x)(v))$ 良定义且 $G \times G$-等变, 且若 $V$ 酉, 其像是 $E_V$.
\end{enumerate}
特别地, 若 $V$ 有限维且酉, 得到满的 $G \times G$-等变映射 $\End(V) \to E_V$.
\end{proposition}

\begin{remark}[非酉情形的矩阵系数]
\label{rem:IV.3.7}
(ii) 提示了把矩阵系数的定义推广到非酉情形的方法: 只把矩阵系数定义为映射 $V^* \otimes_{\CC}V \to \mathcal{C}(G)$ 下纯张量的像.
\end{remark}

\begin{proof}
在此证明中, 我们用 $\cdot$ 记 $G$ 和 $G \times G$ 的所有作用 (这不应引起混淆, 因为每个空间至多有一个作用).
(i). 映射良定义, 因为映射 $V^* \otimes_{\CC}W \to \Hom(V, W)$, 把 $(\Lambda, w)$ 送到 $(v \longmapsto \Lambda(v)w)$ 双线性. 记之为 $\phi$. 为验证它是 $G \times G$-等变的, 只需在纯张量上验证 (因为它们生成 $V^* \otimes_{\CC}W$). 所以设 $\Lambda \in V^*$, $w \in W$, $x, y \in G$. 对每个 $v \in V$,
\[
\phi((x, y) \cdot (\Lambda \otimes w))(v) = \phi((y \cdot \Lambda) \otimes (x \cdot w))(v) = \Lambda(y^{-1} \cdot v)(x \cdot w)
\]
且
\[
((x, y) \cdot \phi(\Lambda \otimes w))(v) = x \cdot (\phi(\Lambda \otimes w)(y^{-1} \cdot v)) = x \cdot (\Lambda(y^{-1} \cdot v)w) = \Lambda(y^{-1} \cdot v)(x \cdot w).
\]
所以
\[
\phi(x \cdot (\Lambda \otimes w)) = x \cdot \phi(\Lambda \otimes w).
\]
假设 $V$ 有限维, 设 $(e_1, \ldots, e_n)$ 是 $V$ 的基, $(e_1^*, \ldots, e_n^*)$ 是对偶基. 定义 $\psi \colon \Hom(V, W) \to V^* \otimes_{\CC}W$, 把 $T$ 送到 $\sum_{i=1}^n e_i^* \otimes T(e_i)$. 让我们证明 $\psi$ 是 $\phi$ 的逆.
若 $j \in \{1, \ldots, m\}$, $w \in W$, 则
\[
\psi(\phi(e_j^* \otimes w)) = \sum_{i=1}^n e_i^* \otimes (\phi(e_j^* \otimes w)(e_i)) = e_j^* \otimes w.
\]
由于元素 $e_j^* \otimes w$ ($j \in \{1, \ldots, n\}$, $w \in W$) 生成 $V^* \otimes_{\CC}W$, 这证明 $\psi \circ \phi = \mathrm{id}$.
反之, 若 $T \in \Hom(V, W)$, 则对每个 $v \in V$,
\[
\phi(\psi(T)) = \sum_{i=1}^n\phi(e_i^* \otimes T(e_i))(v) = \sum_{i=1}^n e_i^*(v)T(e_i) = T(v),
\]
因为 $v = \sum_{i=1}^n e_i^*(v)e_i$. 所以 $\phi(\psi(T)) = T$.
这证明若 $V$ 有限维, 映射 $V^* \otimes_{\CC}W \to \Hom(V, W)$ 是同构. 最后一个结论由此立即得出.
(ii). 映射良定义, 因为映射 $V^* \times V \to \mathcal{C}(G)$, 把 $(\Lambda, v)$ 送到函数 $x \longmapsto \Lambda(\pi(x)(v))$ 双线性. 记之为 $\alpha$. 我们证明 $\alpha$ 是 $G \times G$-等变的. 如前, 只需在纯张量上验证. 所以设 $\Lambda \in V^*$, $v \in V$, $x, y \in G$. 对每个 $z \in G$,
\[
\alpha((x, y) \cdot (\Lambda \otimes v))(z) = \Lambda(x^{-1} \cdot (z \cdot (y \cdot v))) = \Lambda((x^{-1}zy) \cdot v) = ((L_x \circ R_y)(\alpha(\Lambda \otimes v)))(z),
\]
所以 $\alpha((x, y) \cdot (\Lambda \otimes v)) = (L_x \circ R_y)(\alpha(\Lambda \otimes v))$.
最后, 我们证明若 $V$ 酉, $\alpha$ 的像是 $E_V$. 设 $\Lambda \in V^*$. 由于 $V$ 是希尔伯特空间, 存在唯一的 $v \in V$ 使得 $\Lambda = \langle \cdot, v \rangle$. 所以对每个 $w \in V$ 和每个 $x \in G$,
\[
\alpha(\Lambda \otimes w)(x) = \langle \pi(x)(w), v \rangle.
\]
这证明 $\alpha(\Lambda \otimes w)$ 是 $\pi$ 的矩阵系数, 也证明我们以这种方式得到 $\pi$ 的所有矩阵系数.
\end{proof}

现在我们证明只对紧群成立的更强结果. 若 $G$ 是紧群, 固定 $G$ 上的规范化 Haar 测度, 记 $L^p(G)$ 为此测度的 $L^p$ 空间. 注意对每个 $p$ 有 $\mathcal{C}(G) \subset L^p(G)$.

\begin{theorem}[Schur 正交性]
\label{thm:IV.3.8}
设 $G$ 是紧群, $(\pi, V)$ 是 $G$ 的不可约酉表示. 回忆 $V$ 有限维 (由习题 \ref{ex:I.5.5.9}).
\begin{enumerate}
\item[(i)] (Schur 正交性 (Schur orthogonality)) 若 $(\pi', V')$ 是 $G$ 的另一个不等价于 $(\pi, V)$ 的不可约酉表示, 则 $E_{\pi}$ 和 $E_{\pi'}$ 作为 $L^2(G)$ 的子空间正交.
\item[(ii)] 有 $\dim(E_{\pi}) = (\dim V)^2$. 更精确地, 若 $(e_1, \ldots, e_d)$ 是 $V$ 的标准正交基, 记 $\phi_{ij}$ 为函数 $G \to \CC$, $x \longmapsto \langle \pi(x)(e_j), e_i \rangle$, 则 $\{\sqrt{d}\phi_{ij}, 1 \leq i, j \leq d\}$ 是 $E_{\pi}$ 对 $L^2$ 内积的标准正交基.
\item[(iii)] 上面定义的 $G \times G$-等变映射 $\End(V) \to E_{\pi}$ 是同构.
\end{enumerate}
\end{theorem}

\begin{proof}
注意 (iii) 由 (ii) 立即得出, 因为 $\End(V) \to E_{\pi}$ 满射且 (ii) 说 $\dim(E_{\pi}) = (\dim V)^2 = \dim(\End(V))$.
我们证明 (i) 和 (ii). 设 $(\pi', V')$ 是 $G$ 的不可约酉表示, 可以等于 $(\pi, V)$. 若 $A \in \Hom(V, V')$, 定义 $\tilde{A} \in \Hom(V, V')$ 为
\[
\tilde{A} = \int_G \pi'(x)^{-1} \circ A \circ \pi(x)\,\dd x
\]
(注意积分没有问题, 因为表示有限维). 则对每个 $y \in G$,
\begin{align*}
\tilde{A} \circ \pi(y) &= \int_G \pi'(x)^{-1} \circ A \circ \pi(xy)\,\dd x \\
&= \int_G \pi'(xy^{-1})^{-1} \circ A \circ \pi(x)\,\dd x \quad \text{($dx$ 的右不变性)} \\
&= \pi'(y) \circ \tilde{A}.
\end{align*}
换句话说, $\tilde{A}$ 是 $G$-等变的.
设 $v \in V$, $v' \in V'$, 定义 $A \in \Hom(V, V')$ 为 $A(u) = \langle u, v \rangle v'$. 则对所有 $u \in V$ 和 $u' \in V'$,
\begin{align*}
\langle \tilde{A}(u), u' \rangle &= \int_G \langle (\pi'(x)^{-1} \circ A \circ \pi(x))(u), u' \rangle\,\dd x \\
&= \int_G \langle \langle \pi(x)(u), v \rangle\pi'(x)^{-1}(v'), u' \rangle\,\dd x \\
&= \int_G \overline{\langle \pi(x)(u), v \rangle}\langle \pi'(x)(u'), v' \rangle\,\dd x.
\end{align*}
假设 $\pi$ 和 $\pi'$ 不等价. 则由 Schur 引理, 对每个 $A \in \Hom(V, V')$ 有 $\tilde{A} = 0$, 所以由上面的计算, 对所有 $u, v \in V$ 和 $u', v' \in V'$,
\[
\int_G \overline{\langle \pi(x)(u), v \rangle}\langle \pi'(x)(u'), v' \rangle\,\dd x = 0.
\]
这证明 (i).
假设 $\pi = \pi'$, 用 (ii) 的记号. 取 $v = e_i$, $v' = e_{i'}$ ($i, i' \in \{1, \ldots, d\}$), 如前定义 $A$. 再次由 Schur 引理, 存在 $c \in \CC$ 使得 $\tilde{A} = c\,\mathrm{id}_V$. 所以取 $u = e_j$, $u' = e_{j'}$, 由上面的计算得到
\[
\langle \phi_{i,j}, \phi_{i',j'} \rangle_{L^2(G)} = \langle ce_j, e_{j'} \rangle = \begin{cases} c & \text{若 } j = j', \\ 0 & \text{否则}. \end{cases}
\]
另一方面,
\[
cd = \operatorname{Tr}(\tilde{A}) = \int_G \operatorname{Tr}(\pi(x)^{-1} \circ A \circ \pi(x))\,\dd x = \int_G \operatorname{Tr}(A)\,\dd x = \operatorname{Tr}(A).
\]
由于 $A$ 由 $A(w) = \langle w, e_i \rangle e_{i'}$ 定义, 若 $i \neq i'$ 则 $\operatorname{Tr}(A) = 0$, 若 $i = i'$ 则 $\operatorname{Tr}(A) = 1$. 这完成 $\{\sqrt{d}\phi_{ij}, 1 \leq i, j \leq d\}$ 是 $E_{\pi}$ 的标准正交基的证明.
\end{proof}

\section{Peter-Weyl 定理}
设 $G$ 是紧群. 我们把 $L^2(G)$ 看作 $G \times G$ 的表示, 让 $(x, y) \in G \times G$ 按 $L_x \circ R_y = R_y \circ L_x$ 作用. 这限制到第一个 (相应地第二个) 因子是 $G$ 的左 (相应地右) 正则表示.

\begin{definition}[矩阵系数的空间 $E$]
\label{def:IV.4.1}
设 $G$ 是紧群. 记 $E$ 为 $\mathcal{C}(G)$ 中矩阵系数的子空间.
由定理 \ref{thm:IV.3.8}, $E = \bigoplus_{\pi \in \widehat{G}}E_{\pi}$.
\end{definition}

\begin{theorem}[Peter-Weyl 定理: $E$ 稠密]
\label{thm:IV.4.2}
若 $G$ 紧, 则 $E$ 是 $\mathcal{C}(G)$ 的稠密子代数. (对通常的逐点乘法和范数 $\|\cdot\|_{\infty}$.)
\end{theorem}

\begin{proof}
让我们证明 $E$ 被乘法稳定. 注意由命题 \ref{prop:IV.3.2}(iii) 和定理 \ref{thm:IV.2.1}, 对 $G$ 的每个有限维酉表示 $\pi$, 有 $E_{\pi} \subset E$. 设 $(\pi, V)$ 和 $(\pi', V')$ 是 $G$ 的不可约酉表示, $v, w \in V$, $v', w' \in V'$. 回忆我们在习题 \ref{ex:III.6.2.1} 中定义了 $G$ 在 $V \otimes_{\CC}V'$ 上的作用 $\pi \otimes \pi'$\footnote{这只是上面定义的 $G \times G$ 的作用限制到 $G \times G$ 的对角线上.} 以及 $V \otimes_{\CC}V'$ 上的内积.\footnote{这里我们不需要完备化张量积, 因为 $V$ 和 $V'$ 有限维.} 由这些定义, 对每个 $x \in G$,
\[
\langle (\pi \otimes \pi')(x)(v \otimes w), v' \otimes w' \rangle = \langle \pi(x)(v), w \rangle\langle \pi'(x)(v'), w' \rangle.
\]
这证明 $\pi$ 的矩阵系数与 $\pi'$ 的矩阵系数的乘积是 $\pi \otimes \pi'$ 的矩阵系数. 由上面的观察, $\pi \otimes \pi'$ 的每个矩阵系数都在 $E$ 中, 我们完成了.
现在我们证明 $E$ 在 $\mathcal{C}(G)$ 中稠密. 我们已经证明 $E$ 是子代数, 它由命题 \ref{prop:IV.3.2}(iv) 被复共轭稳定, 它包含常数 (它们是 $G$ 在 $\CC$ 上的平凡表示的矩阵系数), 且它由 Gelfand-Raikov 定理 (定理 \ref{thm:III.5.1}) 分离 $G$ 上的点. 所以由 Stone-Weierstrass 定理它在 $\mathcal{C}(G)$ 中稠密.
\end{proof}

\begin{corollary}[$L^p$ 中的稠密性]
\label{cor:IV.4.3}
对每个 $p \in [1, +\infty)$, 子空间 $E$ 在 $L^p(G)$ 中对 $L^p$ 范数稠密.
特别地, 有典则 $G \times G$-等变同构
\[
L^2(G) = \bigoplus_{(\pi, V) \in \widehat{G}}\End(V).
\]
\end{corollary}

最后一个结论通常称为 Peter-Weyl 定理. 它蕴含 $G$ 的左、右正则表示都同构于 $\bigoplus_{\pi \in \widehat{G}}\pi^{\oplus \dim(\pi)}$ 的完备化.

\begin{remark}[经典证明]
\label{rem:IV.4.4}
Peter-Weyl 定理实际上早于 Gelfand-Raikov 定理, 且原始证明用算子 $f * \cdot$ ($f \in L^2(G)$) 在 $L^2(G)$ 上紧这一事实.
\end{remark}

\section{\index{特征标}特征标}

\begin{definition}[特征标 (character)]
\label{def:IV.5.1}
设 $(\pi, V)$ 是拓扑群 $G$ 的连续有限维表示. $\pi$ 的特征标 (character) 是连续映射 $\chi_V = \chi_{\pi} \colon G \to \CC$, $x \longmapsto \operatorname{Tr}(\pi(x))$.
\end{definition}

\begin{remark}[特征标与矩阵系数]
\label{rem:IV.5.2}
若 $(\pi, V)$ 是 $G$ 的有限维表示, $(e_1, \ldots, e_n)$ 是 $V$ 的标准正交基, 则对每个 $x \in G$,
\[
\chi_{\pi}(x) = \sum_{i=1}^n\langle \pi(x)(e_i), e_i \rangle.
\]
所以 $\chi_{\pi} \in E_{\pi}$.
\end{remark}

\begin{definition}[\index{中心函数}中心函数 (central function)]
\label{def:IV.5.3}
若对所有 $x, y \in G$ 有 $f(xyx^{-1}) = f(y)$, 则称函数 $f \colon G \to \CC$ 中心 (central) 或\index{类函数}类函数 (class function).
\end{definition}

这些函数被称为中心的, 因为它们是卷积乘积的中心元, 我们将在第 IV.7 节看到.

\begin{proposition}[特征标的性质]
\label{prop:IV.5.4}
设 $G$ 是拓扑群, $(\pi, V)$ 和 $(\pi', V')$ 是 $G$ 的连续有限维表示. 则:
\begin{enumerate}
\item[(i)] $\chi_{\pi}$ 是中心函数, 且只依赖 $\pi$ 的等价类.
\item[(ii)] $\chi_{V \oplus V'} = \chi_V + \chi_{V'}$.
\item[(iii)] 对每个 $x \in G$, $\chi_{V^*}(x) = \chi_V(x^{-1})$.
\item[(iv)] 对所有 $x, y \in G$,
\[
\chi_{V \otimes_{\CC}V'}(x, y) = \chi_V(x)\chi_{V'}(y) \quad \text{且} \quad \chi_{\Hom(V,V')}(x, y) = \chi_V(x^{-1})\chi_{V'}(y).
\]
\item[(v)] 若 $(\pi, V)$ 可酉化 (例如若 $G$ 紧), 则对每个 $x \in G$ 有 $\chi_V(x^{-1}) = \overline{\chi_V(x)}$.
\end{enumerate}
\end{proposition}

\begin{proof}
(i) 由迹的性质, 即对所有 $A, B \in M_n(\CC)$ 有 $\operatorname{Tr}(AB) = \operatorname{Tr}(BA)$ 得出.
在 $V$ 和 $V'$ 上放任意 Hermitian 内积. 设 $(e_1, \ldots, e_n)$ (相应地 $(e'_1, \ldots, e'_m)$) 是 $V$ (相应地 $V'$) 的标准正交基. 则 $(e_1, \ldots, e_n, e'_1, \ldots, e'_m)$ 是 $V \oplus V'$ 的标准正交基, 所以对每个 $x \in G$,
\[
\chi_{V \oplus V'}(x) = \sum_{i=1}^n\langle \pi(x)(e_i), e_i \rangle + \sum_{j=1}^m\langle \pi'(x)(e'_j), e'_j \rangle = \chi_V(x) + \chi_{V'}(x).
\]
这证明 (ii).
设 $(e_1^*, \ldots, e_n^*)$ 是 $(e_1, \ldots, e_n)$ 的对偶基. 设 $x \in G$. 则若 $A$ 是 $\pi(x^{-1})$ 在基 $(e_1, \ldots, e_n)$ 中的矩阵, $V^*$ 的自同态 $\Lambda \longmapsto \Lambda \circ \pi(x^{-1})$ 在基 $(e_1^*, \ldots, e_n^*)$ 中的矩阵是 $A^T$, 且有
\[
\chi_{V^*}(x) = \operatorname{Tr}(A^T) = \operatorname{Tr}(A) = \chi_V(x^{-1}).
\]
这证明 (iii).
我们证明 $\chi_{V \otimes_{\CC}V'}$ 的公式. 我们在习题 \ref{ex:III.6.2.1} 中见过如何在 $V \otimes_{\CC}V'$ 上放使 $(e_i \otimes e'_j)_{1 \leq i \leq n, 1 \leq j \leq m}$ 成为标准正交基的内积. 所以对所有 $x, y \in G$,
\begin{align*}
\chi_{V \otimes_{\CC}V'}(x, y) &= \sum_{i=1}^n\sum_{j=1}^m\langle \pi(x)(e_i) \otimes \pi'(y)(e'_j), e_i \otimes e'_j \rangle \\
&= \sum_{i=1}^n\sum_{j=1}^m\langle \pi(x)(e_i), e_i \rangle\langle \pi'(y)(e'_j), e'_j \rangle \\
&= \chi_V(x)\chi_{V'}(y).
\end{align*}
现在 $\chi_{\Hom(V,V')}$ 的公式由它、由 (iii) 和命题 \ref{prop:IV.3.6}(i) 得出.
最后, 我们证明 (v). 若 $V$ 可酉化, 可以选择 $V$ 上被 $G$ 不变的 Hermitian 内积. 则对每个 $x \in G$,
\[
\chi_V(x^{-1}) = \sum_{i=1}^n\langle \pi(x^{-1})(e_i), e_i \rangle = \sum_{i=1}^n\langle e_i, \pi(x)(e_i) \rangle = \sum_{i=1}^n\overline{\langle \pi(x)(e_i), e_i \rangle} = \overline{\chi_V(x)}.
\]
\end{proof}

\begin{notation}[不动点子空间]
\label{not:IV.5.5}
若 $(\pi, V)$ 是拓扑群 $G$ 的表示 (连续或不连续), 记
\[
V^G = \{v \in V \mid \forall x \in G,\, \pi(x)(v) = v\}.
\]
这是 $V$ 的闭 $G$-不变子空间.
\end{notation}

\begin{example}[交错算子作为不动点]
\label{exa:IV.5.6}
若 $V$ 和 $W$ 是 $G$ 的两个表示, 则
\[
\Hom(V, W)^G = \Hom_G(V, W).
\]
\end{example}

\begin{theorem}[特征标的积分]
\label{thm:IV.5.7}
设 $G$ 是紧群, $(\pi, V)$ 是 $G$ 的有限维连续表示. 则
\[
\int_G \chi_V(x)\,\dd x = \dim(V^G).
\]
\end{theorem}

\begin{proof}
由于 $V$ 有限维, 可以找到 $V$ 的不可约子表示的有限族 $(V_i)_{i \in I}$ 使得 $V = \bigoplus_{i \in I}V_i$. (参见推论 \ref{cor:I.3.2.9}.) 由命题 \ref{prop:IV.5.4} 有 $\chi_V = \sum_{i \in I}\chi_{V_i}$, 且 $V^G = \bigoplus_{i \in I}V_i^G$. 所以只需对 $V$ 不可约证明定理.
假设 $V$ 是 $G$ 的不可约表示. 由于 $V^G$ 是 $V$ 的 $G$-不变子空间, 有 $V^G = V$ 或 $V^G = \{0\}$. 若 $V^G = V$, 则 $G$ 在 $V$ 上平凡作用, 所以 $V$ 的每个线性子空间都被 $G$ 不变, 故必有 $\dim V = 1$. 另一方面, 对每个 $x \in G$ 有 $\chi_V(x) = \operatorname{Tr}(1) = 1$, 所以 $\int_G\chi_V(x)\,\dd x = 1$. 假设 $V$ 不可约且 $V^G = \{0\}$. 设 $\pi_0$ 是 $G$ 在 $\CC$ 上的平凡表示. 则由定理 \ref{thm:IV.3.8}(i), $L^2(G)$ 的子空间 $E_{\pi}$ 和 $E_{\pi_0}$ 正交. 但 $E_{\pi_0}$ 是常值函数的子空间, 且我们上面 (注记 \ref{rem:IV.5.2}) 看到 $\chi_V \in E_{\pi}$. 所以 $\chi_V$ 与常值函数 $1$ 正交, 这正好意味着 $\int_G\chi_V(x)\,\dd x = 0$.
\end{proof}

\begin{corollary}[特征标的内积]
\label{cor:IV.5.8}
设 $G$ 是紧群, $(\pi, V)$ 和 $(\sigma, W)$ 是 $G$ 的两个连续有限维表示.
\begin{enumerate}
\item[(i)] 有 $\langle \chi_W, \chi_V \rangle_{L^2(G)} = \dim_{\CC}(\Hom_G(V, W))$.
\item[(ii)] 若 $V$ 和 $W$ 不可约且不等价, 则 $\langle \chi_V, \chi_W \rangle_{L^2(G)} = 0$.
\item[(iii)] 表示 $V$ 不可约当且仅当 $\|\chi\|_{L^2(G)} = 1$.
\end{enumerate}
\end{corollary}

\begin{proof}
(i). 让 $G$ 作用于 $\Hom(V, W)$: $x \cdot T = \sigma(x) \circ T \circ \pi(x)^{-1}$. 我们知道 (参见命题 \ref{prop:IV.3.4}) $\Hom_G(V, W) = \Hom(V, W)^G$. 把定理应用于表示 $\Hom(V, W)$, 并用命题 \ref{prop:IV.5.4} 的 (iv) 和 (v) 计算此表示的特征标, 得到:
\begin{align*}
\dim_{\CC}(\Hom_G(V, W)) &= \dim_{\CC}(\Hom(V, W)^G) \\
&= \int_G \chi_{\Hom(V,W)}(x)\,\dd x \\
&= \int_G \chi_V(x)\overline{\chi_W(x)}\,\dd x \\
&= \langle \chi_W, \chi_V \rangle_{L^2(G)}.
\end{align*}
(ii). 这由 (i) 和 Schur 引理 (定理 \ref{thm:I.3.4.1}) 得出, 或由 $\chi_V \in E_V$, $\chi_W \in E_W$ 且 $E_V$ 和 $E_W$ 在 $L^2(G)$ 中正交 (见定理 \ref{thm:IV.3.8}) 得出.
(iii). 若 $V$ 不可约, 则 Schur 引理蕴含 $\End_G(V)$ 是 $1$ 维的, 所以由 (i) 有 $\|\chi_V\|_{L^2(G)} = 1$. 反之, 假设 $\|\chi_V\|_{L^2(G)} = 1$. 把 $V$ 写成 $V = \bigoplus_{i \in I}V_i$, 其中 $I$ 有限, $V_i$ 是 $V$ 的不可约子表示. 由 (ii), $G$ 的不同构不可约表示的特征标在 $L^2(G)$ 中正交, 所以有
\[
\|\chi_V\|_{L^2(G)}^2 = \sum_{W \in \widehat{G}}n_W\|\chi_W\|_{L^2(G)}^2 = \sum_{W \in \widehat{G}}n_W,
\]
其中对每个 $W \in \widehat{G}$,
\[
n_W = |\{i \in I \mid V_i \simeq W\}|.
\]
由于 $\|\chi_V\|_{L^2(G)} = 1$, 存在唯一的 $W \in \widehat{G}$ 使得 $n_W \neq 0$, 且必有 $n_W = 1$. 由整数 $n_W$ 的定义, 这意味着 $V \simeq W$, 所以 $V$ 不可约.
\end{proof}

\begin{corollary}[特征标的线性无关性]
\label{cor:IV.5.9}
设 $G$ 是紧群. 则族 $(\chi_V)_{V \in \widehat{G}}$ (作为 $L^2(G)$ 或 $\mathcal{C}(G)$ 的元素) 线性无关.
\end{corollary}

\begin{proof}
这由前一推论的 (ii) 得出.
\end{proof}

\begin{corollary}[特征标确定表示]
\label{cor:IV.5.10}
设 $\pi$ 和 $\pi'$ 是紧群 $G$ 的两个连续有限维表示. 则 $\pi$ 和 $\pi'$ 等价当且仅当 $\chi_{\pi} = \chi_{\pi'}$.
\end{corollary}

\begin{proof}
若 $\pi$ 和 $\pi'$ 等价, 已知 $\chi_{\pi} = \chi_{\pi'}$. 反之, 假设 $\chi_{\pi} = \chi_{\pi'}$. 把 $\pi$ 和 $\pi'$ 分解成不可约表示的直和:
\[
\pi \simeq \bigoplus_{\rho \in \widehat{G}}\rho^{n_{\rho}}
\]
且
\[
\pi' \simeq \bigoplus_{\rho \in \widehat{G}}\rho^{m_{\rho}},
\]
其中 $n_{\rho}, m_{\rho} \in \Z_{\geq 0}$, 且除有限多个 $\rho \in \widehat{G}$ 外 $n_{\rho} = m_{\rho} = 0$. 由推论 \ref{cor:IV.5.8}, 有 $\chi_{\pi} = \sum_{\rho \in \widehat{G}}n_{\rho}\chi_{\rho}$ 和 $\chi_{\pi'} = \sum_{\rho \in \widehat{G}}m_{\rho}\chi_{\rho}$ (且这是有限和). 由 $\chi_{\rho}$ 的线性无关性, 等式 $\chi_{\pi} = \chi_{\pi'}$ 蕴含对每个 $\rho \in \widehat{G}$ 有 $n_{\rho} = m_{\rho}$, 这进而蕴含 $\pi$ 和 $\pi'$ 等价.
\end{proof}

\section{\index{傅里叶变换}傅里叶变换}

我们仍假设 $G$ 是紧群.
由命题 \ref{prop:I.4.3.4} 和 \ref{prop:I.4.1.3}, 空间 $L^2(G)$ 对卷积乘积实际上是 Banach 代数. 本节回答``如何在 Peter-Weyl 定理给出的分解中看到代数结构?''这个问题.

\begin{definition}[傅里叶变换 (Fourier transform)]
\label{def:IV.6.1}
设 $f \in L^2(G)$. 对每个 $(\pi, V) \in \widehat{G}$, $f$ 在 $\pi$ 处的傅里叶变换是 $V$ 的自同态
\[
\hat{f}(\pi) = \int_G f(x)\pi(x^{-1})\,\dd x = \int_G f(x)\pi(x)^*\,\dd x.
\]
\end{definition}

这显然是 $V$ 的 $\CC$-线性自同态.

\begin{example}[$G = S^1$ 的情形]
\label{exa:IV.6.2}
假设 $G = S^1$. 则我们在习题 \ref{ex:I.5.4.1} 中见过 $\widehat{G} \simeq \Z$, 其中 $n \in \Z$ 对应于表示 $G \to \CC^{\times}$, $e^{2\pi it} \longmapsto e^{2\pi int}$ ($t \in \R$). 所以若 $f \in L^1(G)$, 其傅里叶变换是函数 $\hat{f} \colon \Z \to \CC$, 把 $n$ 送到
\[
\hat{f}(n) = \int_0^1 f(e^{2\pi it})e^{-2\pi int}\,\dd t.
\]
\end{example}

\begin{theorem}[傅里叶变换的性质]
\label{thm:IV.6.3}
\begin{enumerate}
\item[(i)] 对每个 $\pi \in \widehat{G}$, 映射 $L^2(G) \to \End(\pi)$, $f \longmapsto \hat{f}(\pi)$ 是从 $L^2(G)$ 到 $\End(\pi)$ 的相反代数的 $G \times G$-等变 $*$-同态. (注意 $L^2(G) \subset L^1(G)$, 因为 $G$ 紧. 例 \ref{exa:I.4.2.2}(b) 中定义的 $L^1(G)$ 的对合限制为 $L^2(G)$ 上的对合.)
换句话说, 对 $f, g \in L^2(G)$ 和 $x \in G$:
\begin{align*}
\widehat{f * g}(\pi) &= \hat{g}(\pi) \circ \hat{f}(\pi), \\
\widehat{f^*}(\pi) &= (\hat{f}(\pi))^*, \\
\widehat{L_xf}(\pi) &= \hat{f}(\pi) \circ \pi(x)^{-1}, \quad \widehat{R_xf}(\pi) = \pi(x) \circ \hat{f}(\pi).
\end{align*}
(与定理 \ref{thm:I.4.2.6}(i) 比较.)
\item[(ii)] 设 $f \in L^2(G)$. 则对每个 $\pi \in \widehat{G}$, 函数 $\dim(\pi)\operatorname{Tr}(\hat{f}(\pi) \circ \pi(\cdot)) \in L^2(G)$ 是 $f$ 在 $E_{\pi}$ 上的正交投影, 且级数
\[
\sum_{\pi \in \widehat{G}}\dim(\pi)\operatorname{Tr}(\hat{f}(\pi) \circ \pi(\cdot))
\]
在 $L^2(G)$ 中收敛到 $f$ (Fourier 反演公式 (Fourier inversion formula)).
\item[(iii)] 对每个 $f \in L^2(G)$,
\[
\|f\|_2^2 = \sum_{\pi \in \widehat{G}}\dim(\pi)\operatorname{Tr}(\hat{f}(\pi)^* \circ \hat{f}(\pi))
\]
(Parseval 公式 (Parseval formula)).
\end{enumerate}
\end{theorem}

\begin{example}[$S^1$ 的 Fourier 级数]
\label{exa:IV.6.4}
取 $G = S^1$. 则 (ii) 和 (iii) 说: 对每个 $f \in L^2(S^1)$, 级数 $\sum_{n \in \Z}\hat{f}(n)e^{2\pi int}$ 在 $L^1(S^1)$ 中收敛到 $f$, 且
\[
\int_0^1|f(e^{2\pi it})|^2\,\dd t = \sum_{n \in \Z}|\hat{f}(n)|^2.
\]
\end{example}

\begin{proof}
(i). 有
\begin{align*}
\widehat{f * g}(\pi) &= \int_G (f * g)(x)\pi(x^{-1})\,\dd x \\
&= \int_{G \times G}f(y)g(y^{-1}x)\pi(x^{-1})\,\dd x\,\dd y \\
&= \int_{G \times G}f(y)g(x)\pi(x^{-1}y^{-1})\,\dd x\,\dd y \quad \text{(代换 } x' = y^{-1}x\text{)} \\
&= \hat{g}(\pi) \circ \hat{f}(\pi).
\end{align*}
回忆 $f^*(x) = \overline{f(x^{-1})}$, 因为 $\Delta = 1$. 所以
\begin{align*}
\widehat{f^*}(\pi) &= \int_G \overline{f(x^{-1})}\pi(x)^*\,\dd x \\
&= \int_G \overline{f(x)}\pi(x^{-1})^*\,\dd x \\
&= (\hat{f}(\pi))^*.
\end{align*}
最后,
\begin{align*}
\widehat{L_xf}(\pi) &= \int_G f(x^{-1}y)\pi(y^{-1})\,\dd y \\
&= \int_G f(y)\pi(y^{-1}x^{-1})\,\dd y \\
&= \hat{f}(\pi) \circ \pi(x^{-1})
\end{align*}
且
\begin{align*}
\widehat{R_xf}(\pi) &= \int_G f(yx)\pi(y^{-1})\,\dd y \\
&= \int_G f(y)\pi(xy^{-1})\,\dd y \\
&= \pi(x) \circ \hat{f}(\pi).
\end{align*}
(ii). 只需证明第一个结论 (第二个将由 Peter-Weyl 定理得出). 设 $(\pi, V) \in \widehat{G}$. 像定理 \ref{thm:IV.3.8} 中那样, 固定 $V$ 的标准正交基 $(e_1, \ldots, e_d)$, 记 $\phi_{ij}$ 为函数 $G \to \CC$, $x \longmapsto \langle \pi(x)(e_j), e_i \rangle$. 则我们 (在定理 \ref{thm:IV.3.8} 的 (ii) 中) 见过 $\{\sqrt{d}\phi_{ij}, 1 \leq i, j \leq d\}$ 是 $E_{\pi}$ 对 $L^2$ 内积的标准正交基. 所以 $f$ 在 $E_{\pi}$ 上的正交投影是
\[
d\sum_{i,j=1}^d\langle f, \phi_{ij} \rangle_{L^2(G)}\phi_{ij}.
\]
对所有 $i, j \in \{1, \ldots, d\}$,
\begin{align*}
\langle f, \phi_{i,j} \rangle_{L^2(G)} &= \int_G f(x)\overline{\langle e_i, \pi(x)(e_j) \rangle}\,\dd x \\
&= \int_G f(x)\langle \pi(x)^*(e_i), e_j \rangle\,\dd x \\
&= \langle \hat{f}(\pi)(e_i), e_j \rangle.
\end{align*}
设 $y \in G$, 记 $(\hat{f}(\pi)_{i,j})$ 和 $(\pi(y)_{i,j})$ 为 $\hat{f}(\pi)$ 和 $\pi(y)$ 在基 $(e_1, \ldots, e_d)$ 中的矩阵. 则
\[
\hat{f}(\pi)_{i,j} = \langle \hat{f}(\pi)(e_j), e_i \rangle = \langle f, \phi_{j,i} \rangle_{L^2(G)}
\]
且
\[
\pi(y)_{i,j} = \langle \pi(y)(e_j), e_i \rangle = \phi_{ij}(y),
\]
所以
\[
\operatorname{Tr}(\hat{f}(\pi) \circ \pi(y)) = \sum_{i,j=1}^d\hat{f}(\pi)_{j,i}\pi(y)_{i,j} = \sum_{i,j=1}^d\langle \hat{f}(\pi), \phi_{i,j} \rangle_{L^2(G)}\phi_{i,j}(y).
\]
这给出 $f$ 在 $E_{\pi}$ 上的正交投影的所要公式.
(iii). 设 $\pi \in \widehat{G}$, 用 (ii) 证明中的记号. 设 $g = d\operatorname{Tr}(\hat{f}(\pi)^* \circ \hat{f}(\pi))$. 只需证明 $\|g\|_2^2 = d\operatorname{Tr}(\hat{f}(\pi)^* \circ \hat{f}(\pi))$ (因为不同构 $\pi$ 的 $E_{\pi}$ 正交, 由定理 \ref{thm:IV.3.8}). 有
\[
\operatorname{Tr}(\hat{f}(\pi)^* \circ \hat{f}(\pi)) = \sum_{i,j=1}^d|\hat{f}(\pi)_{i,j}|^2 = \sum_{i,j=1}^d|\langle f, \phi_{i,j} \rangle_{L^2(G)}|^2.
\]
另一方面, 由于 $g = d\sum_{i,j=1}^d\langle f, \phi_{ij} \rangle_{L^2(G)}\phi_{ij}$, 得到
\[
\|g\|_{L^2(G)}^2 = d^2\sum_{i,j=1}^d|\langle f, \phi_{i,j} \rangle_{L^2(G)}|^2 = d \cdot d\operatorname{Tr}(\hat{f}(\pi)^* \circ \hat{f}(\pi)).
\]
\end{proof}

\section{特征标与傅里叶变换}
为完成本章, 我们把特征标与傅里叶变换联系起来, 并解释``中心函数''这个名字.
\begin{proposition}[特征标与卷积]
\label{prop:IV.7.1}
设 $f \in L^2(G)$. 则对每个 $x \in G$,
\[
\operatorname{Tr}(\hat{f}(\pi) \circ \pi(x)) = f * \chi_{\pi}(x).
\]
\end{proposition}

由定理 \ref{thm:IV.6.3}, 这说 $f$ 在 $E_{\pi}$ 上的正交投影是 $\dim(\pi)f * \chi_{\pi}$, 所以在 $L^2(G)$ 中有
\[
f = \sum_{\pi \in \widehat{G}}\dim(\pi)f * \chi_{\pi}.
\]

\begin{proof}
有
\begin{align*}
\operatorname{Tr}(\hat{f}(\pi) \circ \pi(x)) &= \int_G f(y)\operatorname{Tr}(\pi(y)^{-1}\pi(x))\,\dd y \\
&= \int_G f(y)\chi_{\pi}(y^{-1}x)\,\dd y \\
&= f * \chi_{\pi}(x).
\end{align*}
\end{proof}

\begin{corollary}[特征标的卷积]
\label{cor:IV.7.2}
对所有 $\pi, \pi' \in \widehat{G}$,
\[
\chi_{\pi} * \chi_{\pi'} = \begin{cases} \dim(\pi)^{-1}\chi_{\pi} & \text{若 } \pi \simeq \pi', \\ 0 & \text{否则}. \end{cases}
\]
\end{corollary}

\begin{proof}
我们知道对每个 $\pi \in \widehat{G}$ 有 $\chi_{\pi} \in E_{\pi}$, $\pi \not\simeq \pi'$ 时 $E_{\pi}$ 和 $E_{\pi'}$ 正交, 且命题说 $\dim(\pi)\chi_{\pi} * \chi_{\pi'}$ 是 $\chi_{\pi}$ 在 $E_{\pi'}$ 上的正交投影. 这立即蕴含推论的公式.
\end{proof}

\begin{definition}[中心函数空间]
\label{def:IV.7.3}
对 $1 \leq p < +\infty$, 记 $ZL^p(G)$ 为 $L^p(G)$ 中中心函数的子空间. 又记 $Z\mathcal{C}(G)$ 为 $\mathcal{C}(G)$ 中中心函数的子空间.
\end{definition}

\begin{proposition}[$L^p$ 的中心]
\label{prop:IV.7.4}
空间 $L^p(G)$ ($1 \leq p < +\infty$) (相应地 $\mathcal{C}(G)$) 对卷积乘积是 Banach 代数, 且 $ZL^p(G)$ (相应地 $Z\mathcal{C}(G)$) 是它的中心.
\end{proposition}

\begin{proof}
设 $p \in [1, +\infty)$, $q \in [1, +\infty)$ 满足 $p^{-1} + q^{-1} = 1$. 由于 $G$ 紧, 常值函数 $1$ 在 $L^q(G)$ 中且 $L^q$ 范数为 $1$, 所以由 H\"older 不等式, $f = f \cdot 1$ 在 $L^1(G)$ 中, 且 $\|f\|_1 \leq \|f\|_p$. 现在推论 \ref{cor:I.4.3.2} 说: 对每个 $g \in L^p(G)$, 函数 $f * g$ 存在且在 $L^p(G)$ 中, 且有 $\|f * g\|_p \leq \|f\|_1\|g\|_p \leq \|f\|_p\|g\|_p$. 这证明 $L^p(G)$ 对 $*$ 是 Banach 代数.
我们证明 $\mathcal{C}(G)$ 对 $*$ 也是 Banach 代数. 若 $f, g \in \mathcal{C}(G)$, 则 $f * g$ 显然存在, 且对每个 $x \in G$,
\[
|f * g(x)| \leq \int_G |f(y)||g(y^{-1}x)|\,\dd y \leq \|f\|_{\infty}\|g\|_{\infty}\int_G 1\,\dd y = \|f\|_{\infty}\|g\|_{\infty}.
\]
所以 $\|f * g\|_{\infty} \leq \|f\|_{\infty}\|g\|_{\infty}$.
最后, 我们证明关于中心的结论. 设 $f \in L^p(G)$, 假设对每个 $g \in L^p(G)$ 有 $f * g = g * f$. 则对每个 $x \in G$ 和每个 $g \in L^p(G)$,
\[
\int_G f(xy)g(y^{-1})\,\dd y = \int_G g(y)f(y^{-1}x)\,\dd y = \int_G f(yx)g(y^{-1})\,\dd y.
\]
这当且仅当 $f(xy) = f(yx)$ 在 $G \times G$ 上几乎处处成立. 对 $f \in \mathcal{C}(G)$ 的证明相同.
\end{proof}

\begin{corollary}[$ZL^2(G)$ 的基]
\label{cor:IV.7.5}
族 $(\chi_{\pi})_{\pi \in \widehat{G}}$ 是 $ZL^2(G)$ 的标准正交基.
\end{corollary}

\begin{proof}
我们已经知道 $\chi_{\pi}$ 在 $ZL^2(G)$ 中且两两正交, 所以只需证明正交于所有 $\chi_{\pi}$ 的中心函数必为 $0$. 设 $f \in ZL^2(G)$. 由下面的引理, 对每个 $\pi \in \widehat{G}$ 有 $(\dim\pi)f * \chi_{\pi} = \langle f, \chi_{\pi} \rangle_{L^2(G)}\chi_{\pi}$, 所以若 $f$ 与每个 $\chi_{\pi}$ 正交, 则由命题 \ref{prop:IV.7.1}, 它在所有空间 $E_{\pi}$ 上的投影为 $0$, 从而由定理 \ref{thm:IV.4.2} 有 $f = 0$.
\end{proof}

\begin{lemma}[中心函数的傅里叶变换]
\label{lem:IV.7.6}
若 $f \in ZL^1(G)$, $\pi \in \widehat{G}$, 则 $(\dim\pi)f * \chi_{\pi} = \langle f, \chi_{\pi} \rangle_{L^2(G)}\chi_{\pi}$.
\end{lemma}

\begin{proof}
我们知道由命题 \ref{prop:IV.7.1} 有 $f * \chi_{\pi} = \operatorname{Tr}(\hat{f}(\pi) \circ \pi(\cdot))$. 对每个 $x \in G$,
\begin{align*}
\hat{f}(\pi) \circ \pi(x) &= \int_G f(y)\pi(y^{-1}x)\,\dd y \\
&= \int_G f(xy^{-1})\pi(y)\,\dd y \\
&= \int_G f(y^{-1}x)\pi(y)\,\dd y \quad \text{(因为 } f \text{ 中心)} \\
&= \int_G f(y)\pi(xy^{-1})\,\dd y \\
&= \pi(x) \circ \hat{f}(\pi).
\end{align*}
所以 $\hat{f}(\pi) \in \End(\pi)$ 是 $G$-等变的. 由 Schur 引理, 这蕴含 $\hat{f}(\pi) = c\,\mathrm{id}$, 其中 $c \in \CC$. 取迹给出
\[
c(\dim\pi) = \operatorname{Tr}(\hat{f}(\pi)) = \int_G f(y)\operatorname{Tr}(\pi(y^{-1}))\,\dd y = \langle f, \chi_{\pi} \rangle_{L^2(G)}.
\]
所以
\[
\langle f, \chi_{\pi} \rangle_{L^2(G)}\chi_{\pi} = (\dim\pi)\operatorname{Tr}(\hat{f}(\pi) \circ \pi(\cdot)) = (\dim\pi)f * \chi_{\pi}.
\]
\end{proof}

\begin{remark}[$ZL^p(G)$ 中的稠密性]
\label{rem:IV.7.7}
事实上, 我们甚至可以证明族 $(\chi_{\pi})_{\pi \in \widehat{G}}$ 在 $ZL^p(G)$ (对每个 $p \in [1, +\infty)$) 和 $ZC(G)$ 中张成稠密子空间. (见 {[}11{]} 的命题 5.25.)
\end{remark}

\begin{remark}[有限群的情形]
\label{rem:IV.7.8}
若 $G$ 有限, 则 $L^2(G)$ 是所有从 $G$ 到 $\CC$ 的函数的空间, $ZL^2(G)$ 是在 $G$ 的共轭类上取常值的函数的空间. 所以上面的命题说 $|\widehat{G}|$ 等于 $G$ 的共轭类数, 且 Peter-Weyl 定理说 $|G| = \sum_{\pi \in \widehat{G}}(\dim\pi)^2$.
\end{remark}

\begin{remark}[$ZL^p(G)$ 的谱]
\label{rem:IV.7.9}
我们特别证明了 Banach 代数 $(ZL^p(G), *)$ ($1 \leq p < +\infty$) 和 $(Z\mathcal{C}(G), *)$ 交换. 我们可以问它们的谱是什么. 事实上, 答案非常简单 (见 {[}11{]} 的定理 5.26): 对每个 $\pi \in \widehat{G}$, 公式 $f \longmapsto (\dim\pi)\int_G f\chi_{\pi}\,\dd\mu$ 定义 $ZL^p(G)$ (相应地 $Z\mathcal{C}(G)$) 上的乘法泛函, 且这诱导从离散集 $\widehat{G}$ 到 $ZL^p(G)$ (相应地 $Z\mathcal{C}(G)$) 的谱的同胚.
\end{remark}

\section{Peter-Weyl 定理的经典证明}
在第 IV.4 节中, 我们给出了使用 Gelfand-Raikov 和 Stone-Weierstrass 定理的 Peter-Weyl 定理的证明. 我们现在解释原始证明. 我们固定紧群 $G$. 回忆这蕴含 $L^2(G) \subset L^1(G)$.
由推论 \ref{cor:I.4.3.2}, 若 $f \in L^1(G)$, $g \in L^2(G)$, 则定义 $f * g$ 和 $g * f$ 的积分收敛, 并定义 $L^2(G)$ 的函数, 使得 $\|f * g\|_2 \leq \|f\|_1\|g\|_2$, $\|g * f\|_2 \leq \|f\|_1\|g\|_2$.

\begin{definition}[$L^2(G)$ 上的卷积算子]
\label{def:IV.8.1}
若 $f \in L^1(G)$, 定义 $L^2(G)$ 的连续线性自同态 $L_f$ 和 $R_f$ 为 $L_f(g) = f * g$, $R_f(g) = g * f$.
\end{definition}

事实上, 由习题 \ref{ex:I.5.6.4} (及其对右乘法的显然类似), $L^1(G)$ 在 $L^2(G)$ 上的这些作用正是定理 \ref{thm:I.4.2.6} 中定义的 $G$ 在 $L^2(G)$ 上的左、右正则表示向群代数 $L^1(G)$ 的延拓. 特别地, $\|L_f\|_{\mathrm{op}} \leq \|f\|_1$, $\|R_f\|_{\mathrm{op}} \leq \|f\|_1$, 这也是推论 \ref{cor:I.4.3.2} 已知的.
在例 \ref{exa:I.4.2.2}(b) 中, 我们在 $L^1(G)$ 上定义了对合 $*$; 由于 $G$ 的模函数为 $1$, 此对合把 $f \in L^1(G)$ 送到由 $f^*(x) = \overline{f(x^{-1})}$ 定义的函数 $f^*$. 由定理 \ref{thm:I.4.2.6}, 对每个 $f \in L^1(G)$ 有 $(L_f)^* = L_{f^*}$, $(R_f)^* = R_{f^*}$.

\begin{theorem}[卷积算子紧]
\label{thm:IV.8.2}
对每个 $f \in L^1(G)$, $L^2(G)$ 的自同态 $R_f$ 和 $L_f$ 紧.
\end{theorem}

\begin{remark}[迹类与 Hilbert-Schmidt]
\label{rem:IV.8.3}
我们还可以证明 $L_f$ 和 $R_f$ 是迹类算子 (trace class operator) (对 $f \in L^1(G)$ 迹为 $\int_G f(x)\,\dd x$), 且对 $f \in L^2(G)$ 是 Hilbert-Schmidt 算子 (Hilbert-Schmidt norm 为 $\|f\|_2$), 但我们不需要.
\end{remark}

\begin{lemma}[核算子]
\label{lem:IV.8.4}
对每个 $F \in L^2(G \times G)$, 公式
\[
T_F(h)(x) = \int_G F(x, y)h(y)\,\dd y
\]
定义连续线性算子 $T_F \colon L^2(G) \to L^2(G)$, 且有 $\|T_F\| \leq \|F\|_2$.\footnote{事实上, 可以证明 $T_F$ 是 Hilbert-Schmidt 算子且其 Hilbert-Schmidt 范数等于 $\|F\|_2$.}
\end{lemma}

\begin{proof}
设 $h \in L^2(G)$. 由 Minkowski 不等式和 Cauchy-Schwarz 不等式,
\begin{align*}
\Bigl(\int_G\Bigl|\int_G F(x, y)h(y)\,\dd y\Bigr|^2\,\dd x\Bigr)^{1/2} &\leq \int_G\Bigl(\int_G|F(x, y)|^2|h(y)|^2\,\dd x\Bigr)^{1/2}\,\dd y \\
&\leq \int_G |h(y)|\Bigl(\int_G|F(x, y)|^2\,\dd x\Bigr)^{1/2}\,\dd y \\
&\leq \|h\|_2\Bigl(\int_G\Bigl(\int_G|F(x, y)|^2\,\dd x\Bigr)^{1/2 \times 2}\,\dd y\Bigr)^{1/2} \\
&= \|h\|_2\|F\|_2.
\end{align*}
这证明 $T_F(h)$ 良定义, 在 $L^2(G)$ 中, 且 $\|T_F(h)\|_2 \leq \|h\|_2\|F\|_2$.
\end{proof}

\begin{lemma}[$L^2(G) \otimes L^2(G) \simeq L^2(G \times G)$]
\label{lem:IV.8.5}
若 $f_1, f_2 \colon G \to \CC$, 定义函数 $u(f_1 \otimes f_2) \colon G \times G \to \CC$ 为 $(f_1 \otimes f_2)(x_1, x_2) = f_1(x_1)f_2(x_2)$. 这诱导 $\CC$-线性等距
\[
u \colon L^2(G) \otimes_{\CC}L^2(G) \to L^2(G \times G),
\]
它单射且像稠密, 从而给出等距同构 $L^2(G)\hat{\otimes}_{\CC}L^2(G) \xrightarrow{\sim} L^2(G \times G)$.
\end{lemma}

证明此引理后, 我们直接把 $L^2(G)\hat{\otimes}_{\CC}L^2(G)$ 与 $L^2(G \times G)$ 等同, 并写 $f_1 \otimes f_2$ 代替 $u(f_1 \otimes f_2)$.

\begin{proof}
$\CC$-线性映射 $u$ 的存在性由张量积的性质得出. 若 $f_1, f_2, g_1, g_2 \in L^2(G)$, 则
\begin{align*}
\langle u(f_1 \otimes f_2), u(g_1 \otimes g_2) \rangle_{L^2(G \times G)} &= \int_{G \times G}u(f_1 \otimes f_2)(x_1, x_2)\overline{u(g_1 \otimes g_2)(x_1, x_2)}\,\dd x_1\,\dd x_2 \\
&= \int_{G \times G}f_1(x_1)f_2(x_2)\overline{g_1(x_1)}\overline{g_2(x_2)}\,\dd x_1\,\dd x_2 \\
&= \langle f_1, g_1 \rangle_{L^2(G)}\langle f_2, g_2 \rangle_{L^2(G)}.
\end{align*}
这蕴含 $u$ 等距.
取 $L^2(G)$ 的 Hilbert 基 $(e_i)_{i \in I}$. 则由上面的计算, $L^2(G \times G)$ 的族 $(u(e_i \otimes e_j))_{i,j \in I}$ 标准正交, 且其张成是 $u$ 的像. 所以我们只需证明正交于每个 $u(e_i \otimes e_j)$ 的 $L^2(G \times G)$ 中的函数为 $0$. 设 $F \in L^2(G \times G)$ 是这种函数. 设 $i \in I$, 考虑函数 $f_i \colon G \to \CC$, $f_i(x) = \int_G F(x, y)\overline{e_i(y)}\,\dd y$. 则 $f_i \in L^2(G)$ (由引理 \ref{lem:IV.8.4}), 且 $f_i$ 与所有 $e_j$ 正交 (由 $F$ 的选取). 由于 $(e_j)_{j \in I}$ 是 $L^2(G)$ 的 Hilbert 基, 必有 $f_i = 0$. 这对每个 $i \in I$ 成立, 所以再次用 $(e_i)_{i \in I}$ 是 $L^2(G)$ 的 Hilbert 基这一事实, 看到函数 $F(x, \cdot)$ 对几乎每个 $x$ 几乎处处为零, 即 $F = 0$ ($L^2(G \times G)$ 中).
\end{proof}

\begin{lemma}[核算子紧]
\label{lem:IV.8.6}
对每个 $F \in L^2(G \times G)$, 引理 \ref{lem:IV.8.4} 的算子 $T_F \colon L^2(G) \to L^2(G)$ 紧.
\end{lemma}

\begin{proof}
由引理 \ref{lem:IV.8.5}, 存在 $L^2(G)$ 的函数族 $(g_i)_{i \in I}$ 和 $(h_i)_{i \in I}$ 使得和 $\sum_{i \in I}g_i \otimes h_i$ 在 $L^2(G \times G)$ 中收敛到 $F$.
对 $I$ 的每个有限子集 $J$, 设 $S_J = \sum_{i,j \in J}g_i \otimes h_j \in L^2(G \times G)$, $T_J = T_{S_J}$. 则对每个 $h \in L^2(G)$, 每个 $x \in G$,
\[
(T_Jh)(x) = \sum_{(i,j) \in J^2}\int_G h(y)g_i(x)\overline{h_j(y)}\,\dd y = \sum_{i \in J}\Bigl(\sum_{j \in J}\int_G h(y)\overline{h_j(y)}\,\dd y\Bigr)g_i(x).
\]
换句话说, 对每个 $h \in L^2(G)$, $T_Jh$ 在由 $g_i$ ($i \in J$) 张成的 $L^2(G)$ 的有限维子空间中. 所以算子 $T_J$ 有有限秩.
为证明 $T_F$ 紧, 由习题 \ref{ex:I.5.5.9}(f) 只需证明它是 $J$ 变大时算子 $T_J$ 的极限. 但这由引理 \ref{lem:IV.8.4} 和 $S_J$ 在 $L^2(G \times G)$ 中收敛到 $F$ 的事实得出.
\end{proof}

\begin{proof}[定理 \ref{thm:IV.8.2} 的证明]
我们证明 $L_f$ 的结果; $R_f$ 的证明类似.
首先假设 $f \in L^2(G)$. 考虑函数 $K \colon G \times G \to \CC$, $(x, y) \longmapsto f(xy^{-1})$. 则 $K \in L^2(G \times G)$, 且 $L_f = T_K$. 所以结果由引理 \ref{lem:IV.8.6} 得出.
一般情形, 结果由紧算子的极限紧 (见习题 \ref{ex:I.5.5.9}(f))、$L^2(G)$ 在 $L^1(G)$ 中稠密以及 $f \longmapsto L_f$ (相应地 $f \longmapsto R_f$) 是从 $L^1(G)$ 到 $\End(L^2(G))$ 的连续映射这一事实得出.
\end{proof}

回忆 $\mathcal{C}(G)$ 的子空间 $E$ 在第 IV.4.1 节中定义; 这是 $G$ 的有限维表示的矩阵系数的空间.

\begin{proposition}[有限维子表示]
\label{prop:IV.8.7}
设 $f \in L^2(G)$. 则下列条件等价:
\begin{enumerate}
\item[(i)] $\operatorname{Span}\{L_xf, x \in G\}$ 有限维;
\item[(ii)] $\operatorname{Span}\{R_xf, x \in G\}$ 有限维;
\item[(iii)] $f \in E$.
\end{enumerate}
\end{proposition}

\begin{proof}
我们首先证明 (i) 蕴含 (ii) 和 (iii). 若 $f \in E$, 则 $f$ 是 $\sum_{\pi \in \widehat{G}}f_{\pi}$ 的有限和, 其中 $f_{\pi} \in E_{\pi}$. 所以只需证明对 $\pi \in \widehat{G}$, $E_{\pi}$ 的每个元素在左 (相应地右) 正则表示中生成有限维子空间; 但这由 $E_{\pi}$ 有限维且被算子 $L_x$ (相应地 $R_x$) ($x \in G$) 稳定这一事实立即得出.
现在我们证明 (ii) 蕴含 (i). (证明 (iii) 蕴含 (i) 类似.) 所以设 $f \in L^2(G)$, 假设存在包含 $f$ 且被所有 $L_x$ ($x \in G$) 稳定的 $L^2(G)$ 的有限维子空间 $V$. 我们将证明 $f$ 是 $V^*$ 的矩阵系数. 设 $\langle \cdot, \cdot \rangle_V$ 是 $L^2(G)$ 的内积到 $V$ 的限制, $\rho \colon G \to \GL(V)$ 是 $G$ 在 $V$ 上的作用, $\Pi \colon L^2(G) \to V$ 是正交投影. 假设 $\psi \in \mathcal{C}(G)$ 取实值且对每个 $x \in G$ 有 $\psi(x^{-1}) = \psi(x)$. 特别地, $\psi \in L^2(G)$, 所以由命题 \ref{prop:I.4.3.4}, $f * \psi$ 连续. 对每个 $x \in G$,
\begin{align*}
f * \psi(x) &= L_{x^{-1}}(f * \psi)(1) \\
&= (L_{x^{-1}}f) * \psi(1) \quad \text{(由命题 \ref{prop:I.4.1.3})} \\
&= \int_G (L_{x^{-1}}f)(y)\overline{\psi(y)}\,\dd y \quad \text{(因为 } \psi \text{ 取实值且 } \psi(y^{-1}) = \psi(y)\text{)} \\
&= \langle L_{x^{-1}}f, \psi \rangle_{L^2(G)} \\
&= \langle \rho(x^{-1})\Pi(f), \Pi(\psi) \rangle_V \quad \text{(因为 } L_{x^{-1}}f \in V\text{)} \\
&= \langle \rho(x)\Pi(\psi), \Pi(f) \rangle_V.
\end{align*}
由命题 \ref{prop:IV.3.2}(iv), $f * \psi$ 是 $V^*$ 的矩阵系数, 即 $f * \psi \in E_{V^*}$. 现在我们选取近似单位元 $(\psi_U)_{U \in \mathcal{U}}$ (定义 \ref{def:I.4.1.7}, 由命题 \ref{prop:I.4.1.8} 存在). 由我们刚做的计算, 对每个 $U \in \mathcal{U}$ 有 $f * \psi_U \in E_{V^*}$. 但由推论 \ref{cor:I.4.3.3}, $f$ 是族 $(f * \psi_U)_{U \in \mathcal{U}}$ 在 $L^2(G)$ 中的极限. 由于 $E_{V^*}$ 有限维 (由命题 \ref{prop:IV.3.2}(ii)), 它是 $L^2(G)$ 的闭子空间, 所以 $f$ 也在 $E_{V^*}$ 中.
\end{proof}

\begin{proof}
现在我们解释如何不用 Gelfand-Raikov 定理证明定理 \ref{thm:IV.4.2} (即 $E$ 在 $\mathcal{C}(G)$ 中稠密这一事实). 只需证明 $E$ 在 $L^2(G)$ 中稠密. 考虑 $G$ 在 $L^2(G)$ 上的左正则表示. 由命题 \ref{prop:IV.8.7}, $E$ 包含 $L^2(G)$ 的每个有限维子表示, 事实上它是 $L^2(G)$ 的所有有限维子表示的和. 由于 $L^2(G)$ 是希尔伯特空间, 为证明 $E$ 在 $L^2(G)$ 中稠密, 只需证明 $E^{\perp} = \{0\}$. 所以设 $f \in E^{\perp}$. 选取近似单位元 $(\psi_U)_{U \in \mathcal{U}}$ (定义 \ref{def:I.4.1.7}, 由命题 \ref{prop:I.4.1.8} 存在). 若 $U \in \mathcal{U}$, 则 $\psi_U^* = \psi_U$, 所以 $R_{\psi_U}$ 自伴, 且由定理 \ref{thm:IV.8.2} 它紧. 所以由自伴紧算子的谱定理 (定理 \ref{thm:IV.1.3}), $\Ker(R_{\psi_U})$ 是 $\bigoplus_{\lambda \in \CC^{\times}}\Ker(R_{\psi_U} - \lambda\mathrm{id})$ 的和的闭包的正交补, 且对每个 $\lambda \in \CC^{\times}$, $\Ker(R_{\psi_U} - \lambda\mathrm{id})$ 有限维. 又, $R_{\psi_U}$ 与 $G$ 的左作用交换 (由命题 \ref{prop:I.4.1.3}), 所以每个空间 $\Ker(R_{\psi_U} - \lambda\mathrm{id})$ 都是 $L^2(G)$ 的子表示. 由于由假设 $f$ 与 $L^2(G)$ 的每个有限维子表示正交, 这蕴含 $f \in \Ker(R_{\psi_U})$, 即 $f * \psi_U = 0$. 但 $f$ 是 $(f * \psi_U)_{U \in \mathcal{U}}$ 在 $L^2(G)$ 中的极限 (由推论 \ref{cor:I.4.3.3}), 所以我们得出结论 $f = 0$.
\end{proof}

\section{习题}

\begin{exercise}[$\SU(d)$ 的表示]
\label{ex:IV.9.1}
设 $d$ 是正整数.
\begin{enumerate}
\item[(a)] 设 $T$ 是 $\UU(d)$ 与对角矩阵集的交. 证明
\[
T = \left\{\begin{pmatrix} z_1 & \cdots & 0 \\ \vdots & \ddots & \vdots \\ 0 & \cdots & z_d \end{pmatrix},\, z_1, \ldots, z_d \in S^1\right\}.
\]
\item[(b)] 证明 $\UU(d)$ 的每个元素都在 $\UU(d)$ 中与 $T$ 的元素共轭.
\item[(c)] 证明 $\SU(d)$ 的每个元素都在 $\SU(d)$ 中与 $T_0 := T \cap \SU(d)$ 的元素共轭.
\item[(d)] 证明 $\SU(d)$ 的有限维表示 $V$ 由 $\chi_V|_{T_0}$ 在等价意义下唯一确定.
\end{enumerate}
我们现在取 $d = 2$. 回忆习题 \ref{ex:I.5.5.1} 中定义的 $\SU(2)$ 的不可约表示 $V_n$ ($n \geq 0$).
\begin{enumerate}
\item[(a)] 计算 $\chi_{V_n}$ 到 $T_0$ 的限制.
\item[(b)] 设 $(\rho, V)$ 是 $\SU(2)$ 的有限维表示. 证明存在 $m \geq 1$ 和非负整数 $a_0, \ldots, a_m$ 使得对每个 $z \in S^1$,
\[
\chi_V\left(\begin{pmatrix} z & 0 \\ 0 & z^{-1} \end{pmatrix}\right) = a_0 + \sum_{i=1}^m a_i(z^i + z^{-i}).
\]
\item[(c)] 证明存在整数 $c_n \in \Z$ ($n \geq 0$) 使得对足够大的 $n$ 有 $c_n = 0$, 且
\[
\chi_V = \sum_{n \geq 0}c_n\chi_{V_n}.
\]
\item[(d)] 证明 (c) 的整数 $c_n$ 都非负.
\item[(e)] 若 $V$ 不可约, 证明存在 $n \geq 0$ 使得 $V \simeq V_n$.
\end{enumerate}
\end{exercise}

\begin{solution}
(a). 设 $x = \begin{pmatrix} z_1 & \cdots & 0 \\ \vdots & \ddots & \vdots \\ 0 & \cdots & z_d \end{pmatrix} \in M_d(\CC)$, 其中 $z_1, \ldots, z_d \in \CC$. 则 $x \in T$ 当且仅当 $xx^* = I_d$. 由于 $x^* = \begin{pmatrix} \bar{z}_1 & \cdots & 0 \\ \vdots & \ddots & \vdots \\ 0 & \cdots & \bar{z}_d \end{pmatrix}$, 此条件等价于 $|z_1| = \cdots = |z_d| = 1$.
(b). 设 $x \in \UU(d)$. 则 $x$ 正规, 所以由谱定理, 它可以在 $\CC^n$ 的标准正交基中对角化. 这意味着存在 $y \in \UU(d)$ 使得 $yxy^{-1}$ 对角, 即 $yxy^{-1} \in T$.
(c). 设 $x \in \SU(d)$. 由 (b), 存在 $y \in \UU(d)$ 使得 $yxy^{-1} \in T$. 有 $\det(yxy^{-1}) = \det(x) = 1$, 所以 $yxy^{-1}$ 实际上在 $T_0$ 中. 设 $c = \det(y) \in \CC^{\times}$. 选取 $c' \in \CC$ 使得 $(c')^d = c$; 由于 $|c| = 1$, 也有 $|c'| = 1$. 则 $y' := (c')^{-1}y$ 的行列式为 $1$, 所以在 $\SU(d)$ 中, 且 $y'x(y')^{-1} = yxy^{-1}$.
(d). 设 $V, W$ 是 $\SU(d)$ 的两个有限维表示, 假设 $\chi_V = \chi_W$ 在 $T_0$ 上. 由 (c) 和 $\chi_V$, $\chi_W$ 是中心函数这一事实, 这蕴含 $\chi_V = \chi_W$ 在整个 $\SU(d)$ 上. 但此时 $V$ 和 $W$ 由推论 \ref{cor:IV.5.10} 等价.
(e). 设 $x = \begin{pmatrix} x_1 & 0 \\ 0 & x_2 \end{pmatrix} \in T_0$. 注意 $x_1x_2 = 1$, 所以 $x_2 = x_1^{-1} = \bar{x}_1$. 我们计算 $x$ 在 $V_n$ 的基 $(t_1^kt_2^{n-k})_{0 \leq k \leq n}$ 上的作用. 对 $0 \leq k \leq n$,
\[
x \cdot t_1^kt_2^{n-k} = (x_1^{-1}t_1)^k(x_2^{-1}t_2)^{n-k} = x_1^{n-2k}t_1^kt_2^{n-k}.
\]
所以
\[
\chi_{V_n}(x) = \sum_{k=0}^n x_1^{n-2k}.
\]
(f). 我们把 $S^1$ 通过连续群态射 $z \longmapsto \begin{pmatrix} z & 0 \\ 0 & z^{-1} \end{pmatrix}$ 嵌入 $\SU(2)$. 注意这诱导拓扑群的同构 $S^1 \xrightarrow{\sim} T_0$. 则 $\rho|_{S^1}$ 是 $S^1$ 的有限维表示, 所以它是不可约表示的有限直和. 我们知道 (从习题 \ref{ex:I.5.4.1}) $S^1$ 的每个不可约表示都具有 $\rho_m \colon z \longmapsto z^m$ ($m \in \Z$) 的形式, 所以存在非负整数 $a_m$ ($m \in \Z$) 使得除有限多个 $m$ 外 $a_m = 0$, 且 $\rho|_{S^1} \simeq \bigoplus_{m \in \Z}\rho_m^{a_m}$. 特别地, 对每个 $z \in S^1$,
\[
\chi_V\left(\begin{pmatrix} z & 0 \\ 0 & z^{-1} \end{pmatrix}\right) = \sum_{m \in \Z}a_mz^m.
\]
设 $y = \begin{pmatrix} 0 & i \\ i & 0 \end{pmatrix}$. 则 $y \in \SU(2)$, 且对每个 $z \in S^1$,
\[
y\begin{pmatrix} z & 0 \\ 0 & z^{-1} \end{pmatrix}y^{-1} = \begin{pmatrix} z^{-1} & 0 \\ 0 & z \end{pmatrix}.
\]
由于 $V$ 是 $\SU(2)$ 的表示, 函数 $\chi_V$ 在 $\SU(2)$ 上中心, 所以
\[
\chi_V\begin{pmatrix} z & 0 \\ 0 & z^{-1} \end{pmatrix} = \chi_V\begin{pmatrix} z^{-1} & 0 \\ 0 & z \end{pmatrix}.
\]
这蕴含对每个 $m \in \Z$ 有 $a_{-m} = a_m$, 所以我们得到所要的结论.
(g). 设 $M$ 是可以写成 $\chi(z) = a_0 + \sum_{m \geq 1}a_m(z^m + z^{-m})$ ($a_0, a_1, \ldots \in \Z$ 且对足够大的 $m$ 有 $a_m = 0$) 的函数 $\chi \colon S^1 \to \Z$ 的 $\Z$-模. 由 (f), $\chi_V$ 到 $S^1$ 的限制在 $M$ 中.
$M$ 在 $\Z$ 上的基由函数 $(\chi_0 = 1, \chi_1 = z + z^{-1}, \chi_2 = z^2 + z^{-2}, \ldots)$ 组成. 另一方面, 我们在 (e) 中见过 $\chi_{V_n} = \sum_{0 \leq k \leq n}\chi_k$. 所以表示 $(\chi_{V_n})_{n \geq 0}$ 在基 $(\chi_n)_{n \geq 0}$ 中的 (无穷) 矩阵是上三角的, 对角线上为 $1$, 这意味着它可逆, 即 $(\chi_{V_n})_{n \geq 0}$ 也是 $M$ 在 $\Z$ 上的基. (如果你不喜欢这样, 也很容易从把 $\chi_{V_n}$ 在基 $(\chi_m)_{m \geq 0}$ 中表示的公式对 $n$ 归纳证明 $(\chi_{V_0}, \ldots, \chi_{V_n})$ 线性无关且张成与 $(\chi_0, \ldots, \chi_n)$ 相同的 $\Z$-子模.)
题目的结论由此立即得出.
(h). 我们知道函数 $\chi_{V_n}$ 在 $L^2(\SU(2))$ 中两两正交 (由推论 \ref{cor:IV.5.8}). 所以对每个 $n \geq 0$,
\[
c_n = \langle \chi_V, \chi_{V_n} \rangle_{L^2(\SU(2))}.
\]
由同一个推论, 右端也等于 $\dim_{\CC}(\Hom_{\SU(2)}(V, V_n))$, 它是非负整数.
(i). 若 $V$ 不可约, 则由 (h) 的证明中的最后一个公式 (和 Schur 引理), 除非 $V \simeq V_n$ 否则 $c_n = 0$. 所以, 若不存在 $n \geq 0$ 使得 $V \simeq V_n$, 则有 $\chi_V = 0$, 从而 $V = 0$, 这不可能.
\end{solution}

\begin{exercise}[有限群: 忠实的有限维表示]
\label{ex:IV.9.2}
设 $G$ 是紧群, $(\pi, V)$ 是 $G$ 的忠实有限维连续表示. (回忆这意味着 $\pi \colon G \to \GL(V)$ 单射.) 本题的目标是证明: 若 $G$ 有限, 则 $G$ 的每个不可约表示都是形如 $V^{\otimes n} \otimes (V^*)^{\otimes m}$ (对某个 $n, m \geq 1$) 的表示的直和项, 其中记号 $V^{\otimes n}$ 表示 $\underbrace{V \otimes \cdots \otimes V}_{n \text{ 次}}$, 对 $(V^*)^{\otimes m}$ 类似.
\begin{enumerate}
\item[(a)] 设 $\mathbf{1}$ 是 $G$ 在 $\CC$ 上的平凡表示. 证明只需证明 $G$ 的每个不可约表示都是形如 $(V \oplus V^* \oplus \mathbf{1})^{\otimes N}$ (对某个 $N \geq 1$) 的表示的直和项.
\item[(b)] 设 $W$ 是 $G$ 的不可约表示. 证明 $W$ 是 $(V \oplus V^* \oplus \mathbf{1})^{\otimes N}$ 的直和项当且仅当 $\int_G (1 + 2\Re\chi_V(x))^N\chi_W(x)\,\dd x \neq 0$.
\end{enumerate}
从现在起, 我们假设 $G$ 有限, 固定 $G$ 的有限维表示 $W$, 对每个 $N \in \Z_{\geq 0}$ 写
\[
S_N = \sum_{x \in G}(1 + 2\Re\chi_V(x))^N\chi_W(x).
\]
设 $d = \dim V$.
\begin{enumerate}
\item[(c)] 若 $x \neq 1$, 证明当 $N \to +\infty$ 时 $(1 + 2\Re\chi_V(x))^N\chi_W(x) = o((1 + 2d)^N)$.
\item[(d)] 若 $W \neq 0$, 证明对足够大的 $N$ 有 $S_N \neq 0$.
\end{enumerate}
\end{exercise}

\begin{solution}
(a). 让我们对 $N$ 归纳证明: 对每个 $N \in \Z_{\geq 1}$, 有 $G$-等变同构
\[
(\mathbf{1} \oplus V \oplus V^*)^{\otimes N} \simeq \bigoplus_{\substack{k,l,m \geq 0\\ k+l+m=N}}V^{\otimes k} \otimes (V^*)^{\otimes l}\,^{\oplus \frac{N!}{k!l!m!}}.
\]
这显然蕴含 (a) 的结果.
对 $N = 1$, 它由对 $G$ 的每个表示 $W$ 有 $\mathbf{1} \otimes W \simeq W$ 这一事实得出. 假设对 $N$ 已知结果, 证明对 $N+1$. 有
\begin{align*}
(\mathbf{1} \oplus V \oplus V^*)^{\otimes N+1} &\simeq (\mathbf{1} \oplus V \oplus V^*)^{\otimes N} \otimes (\mathbf{1} \oplus V \oplus V^*) \\
&\simeq (\mathbf{1} \oplus V \oplus V^*) \otimes \bigoplus_{\substack{k,l,m \geq 0\\ k+l+m=N}}V^{\otimes k} \otimes (V^*)^{\otimes l}\,^{\oplus \frac{N!}{k!l!m!}} \\
&\simeq \bigoplus_{\substack{k,l,m \geq 0\\ k+l+m=N}}V^{\otimes k} \otimes (V^*)^{\otimes l}\,^{\oplus \frac{N!}{k!l!m!}} \oplus \bigoplus_{\substack{k,l,m \geq 0\\ k+l+m=N}}V^{\otimes k+1} \otimes (V^*)^{\otimes l}\,^{\oplus \frac{N!}{k!l!m!}} \\
&\qquad \oplus \bigoplus_{\substack{k,l,m \geq 0\\ k+l+m=N}}V^{\otimes k} \otimes (V^*)^{\otimes l+1}\,^{\oplus \frac{N!}{k!l!m!}} \\
&\simeq \bigoplus_{\substack{k,l,m \geq 0\\ k+l+m=N+1}}V^{\otimes k} \otimes (V^*)^{\otimes l}\,^{\oplus \bigl(\frac{N!}{k!l!(m-1)!} + \frac{N!}{(k-1)!l!m!} + \frac{N!}{k!(l-1)!m!}\bigr)} \\
&\simeq \bigoplus_{\substack{k,l,m \geq 0\\ k+l+m=N+1}}V^{\otimes k} \otimes (V^*)^{\otimes l}\,^{\oplus \frac{(N+1)!}{k!l!m!}}.
\end{align*}
(b). 由 $G$ 的有限维表示的半单性 (推论 \ref{cor:I.3.2.9}) 和 Schur 引理 (定理 \ref{thm:I.3.4.1}), 表示 $W$ 是 $(\mathbf{1} \oplus V \oplus V^*)^{\otimes N}$ 的直和项当且仅当 $\Hom_G((\mathbf{1} \oplus V \oplus V^*)^{\otimes N}, W) \neq 0$. 由推论 \ref{cor:IV.5.8}, 这当且仅当 $\langle \chi_{(\mathbf{1} \oplus V \oplus V^*)^{\otimes N}}, \chi_W \rangle_{L^2(G)} \neq 0$. 所以结论由 $\chi_{(\mathbf{1} \oplus V \oplus V^*)^{\otimes N}} = (1 + \chi_V + \bar{\chi}_V)^N = (1 + 2\Re\chi_V)^N$ 得出, 这是命题 \ref{prop:IV.5.4} 的直接推论.
(c). 由于 $G$ 紧, 表示 $(\pi, V)$ 可酉化, 所以可以选择同构 $V \simeq \CC^d$ 使得 $\pi(G) \subset \UU(d)$. 设 $z_1, \ldots, z_d$ 是 $\pi(x)$ 的特征值. 由于 $\pi$ 忠实, $\pi(x) \neq 1$, 所以 $z_i$ 中至少有一个不等于 $1$ (我们用 $\pi(x)$ 可对角化这一事实); 所以可以假设 $z_1 \neq 1$. 由于 $|z_1| = 1$, 这蕴含 $-1 \leq \Re z_1 < 1$, 所以 $1 - 2d \leq 1 + 2\sum_{i=1}^d\Re(z_i) < 1 + 2d$, 且 $1 + 2\sum_{i=1}^d\Re z_i < 1 + 2d$.
最后得到
\[
|(1 + 2\Re\chi_V(x))^N\chi_W(x)| \leq (\dim W)\Bigl|1 + 2\sum_{i=1}^d\Re z_i\Bigr|^N = o((1 + 2d)^N).
\]
(d). 由于 $G$ 有限, (c) 蕴含
\[
\sum_{x \in G-\{1\}}(1 + 2\Re\chi_V(x))^N\chi_W(x) = o((1 + 2d)^N).
\]
另一方面, $(1 + 2\Re\chi_V(1))^N\chi_W(1) = (\dim W)(1 + 2d)^N$. 所以
\[
S_N = (\dim W)(1 + 2d)^N + o((1 + 2d)^N),
\]
这蕴含对足够大的 $N$ 有 $S_N \neq 0$.
\end{solution}

\begin{exercise}[任意紧群: 忠实表示与 $V^{\otimes n} \otimes (V^*)^{\otimes m}$]
\label{ex:IV.9.3}
本题的目标是把习题 \ref{ex:IV.9.2} 推广到任意紧群 $G$, 假设关于 Haar 测度的某些条件. 在下一个习题中, 我们给出用矩阵系数的同一结果的另一种方法.
设 $(\rho, V)$ 是 $G$ 的忠实有限维连续表示. 我们要证明 $G$ 的任何不可约表示都是某个 $V^{\otimes N} \otimes (V^*)^{\otimes M}$ 的直和项. 固定 $G$ 上的规范化 Haar 测度 $\mu$.
\begin{enumerate}
\item[(a)] 证明存在同构 $V \simeq \CC^n$ 使得 $\rho$ 诱导 $G$ 与 $\UU(n)$ 的闭子群之间的 (拓扑群的) 同构.
\end{enumerate}
从现在起, 我们假设 $G$ 是 $\UU(n)$ 的闭子群, $V = \CC^n$, $\rho \colon G \to \GL_n(\CC)$ 是包含. 设 $(\pi, W)$ 是 $G$ 的连续非零有限维表示. 定义 $f \colon \UU(n) \to \CC$ 和 $g \colon G \to \CC$ 为 $f(x) = 1 + \operatorname{Tr}(x) + \overline{\operatorname{Tr}(x)}$, $g(x) = \chi_W(x)$.
像习题 \ref{ex:I.5.5.4} 那样定义
\[
L_0 = \{X \in M_n(\CC) \mid \forall t \in \R,\, e^{tX} \in \UU(n)\}
\]
和
\[
L = \{X \in M_n(\CC) \mid \forall t \in \R,\, e^{tX} \in G\}.
\]
回忆我们在习题 \ref{ex:I.5.5.4} 中证明了: 若 $\Omega$ 是 $L$ 中 $0$ 的足够小的邻域, 则 $\exp$ 诱导 $\Omega$ 与 $\exp(\Omega)$ 之间的同胚, 且 $\exp(\Omega)$ 是 $1$ 在 $G$ 中的邻域. 选取同构 $L_0 \simeq \R^m$, 设 $dX$ 是此同构给出的 $L_0$ 上的 Lebesgue 测度. 我们假设: ($*$) 对足够小的 $\Omega$, 存在 $c \in \R_{>0}$ 使得同胚 $\exp \colon \Omega \xrightarrow{\sim} \exp(\Omega)$ 下 Haar 测度 $\mu$ 的原像具有 $h(X)dX$ 的形式, 其中当 $X \to 0$ 时 $h(X) = c + O(\|X\|)$.
注记. 这总是真的, 但我们没有工具证明它. 事实上, $\GL_n(\CC)$ 的闭子群 $G$ 是 Lie 群, 所以 Haar 测度由 $G$ 上的左不变微分形式 $\omega$ 给出, 即 $\omega_g$ 是 $\omega_1$ 经 $g$ 的左平移的拉回. 这和一番努力给出所要的渐近公式.
用 Weyl 积分公式 (见例如 {[}22{]} 的定理 7.16) 来证明本习题中的所有估计实际上会更简单, 因为我们积分的函数是中心函数. 然而, 这需要更多理论 (特别是极大环面).
\begin{enumerate}
\item[(a)] 证明对每个 $x \in \UU(n)$, 有 $f(x) = 1 + 2\sum_{i=1}^n\cos\theta_i$, 其中 $e^{i\theta_1}, \ldots, e^{i\theta_n}$ 是 $x$ 的特征值.
\item[(b)] 若 $\Omega$ 是 $L_0$ 中 $0$ 的邻域, 证明存在 $\delta > 0$ 使得对每个 $x \notin \exp(\Omega)$ 和每个 $N \geq 1$,
\[
|f(x)^N| \leq (1 + 2n - \delta)^N.
\]
\item[(c)] 若 $\Omega$ 是 $L$ 中 $0$ 的邻域, $U = \exp(\Omega)$, 证明存在 $\delta > 0$ 和 $C \in \R_{>0}$ 使得对每个 $N \geq 1$,
\[
\Bigl|\int_{G-U}f(x)^Ng(x)\,\dd\mu(x)\Bigr| \leq C(1 + 2n - \delta)^N.
\]
\item[(d)] 证明当 $X \to 0$ ($L_0$ 中) 时
\[
f(e^X) = (2n + 1)e^{-K(X) + O(\|X\|^4)},
\]
其中 $K(X) = \frac{1}{1+2n}\|X\|^2 = \frac{1}{1+2n}\operatorname{Tr}(X^*X)$.
\item[(e)] 证明若 $\Omega$ 是 $L$ 中以 $0$ 为中心 (有限半径) 的球, 则存在 $D \in \R_{>0}$ 使得当 $N \to +\infty$ 时
\[
\int_{\Omega}e^{-NK(X)}g(e^X)\,\dd X \sim D \cdot N^{-\frac{1}{2}\dim L}.
\]
(提示: 证明当 $X \to 0$ ($L$ 中) 时 $g(e^X) = \dim W + O(\|X\|)$.)
\item[(f)] 证明存在 $1$ 在 $G$ 中的邻域 $U$ 和 $E \in \R_{>0}$ 使得当 $N \to +\infty$ 时
\[
\int_U f(x)^Ng(x)\,\dd\mu(x) \sim E\frac{(2n + 1)^N}{N^{\frac{1}{2}\dim L}}.
\]
\item[(g)] 证明若 $N$ 足够大则 $\int_G f(x)^Ng(x)\,\dd\mu(x) \neq 0$.
\end{enumerate}
\end{exercise}

\begin{solution}
(a). 由于 $G$ 紧, 表示 $(\rho, V)$ 可酉化. 这意味着存在同构 $V \simeq \CC^n$ 使得 $\rho(G) \subset \UU(n)$. 由于表示 $(\rho, V)$ 忠实, 态射 $\rho$ 单射, 所以 $\rho \colon G \to \UU(n)$ 是单射连续映射. 由于 $G$ 紧, 此映射是到其像的同胚.
(b). 设 $D$ 是对角线元素为 $e^{i\theta_1}, \ldots, e^{i\theta_n}$ 的对角矩阵. 由于 $x$ 与 $x^* = x^{-1}$ 交换, 谱定理蕴含存在 $A \in \UU(n)$ 使得 $D = AxA^{-1}$. 由于 $f$ 显然是 $\UU(n)$ 上的中心函数 (甚至在 $\GL_n(\CC)$ 上), 有 $f(x) = f(D)$. 但
\[
f(e^D) = 1 + \sum_{j=1}^n\Re(e^{i\theta_j}) = 1 + 2\sum_{j=1}^n\cos(\theta_j).
\]
(c). 由 (b), 对每个 $x \in \UU(n)$, $f(x) \in \R$ 且 $1 - 2n \leq f(x) \leq 1 + 2n$. 而且, 等式 $f(x) = 1 + 2n$ 只在 $x$ 的所有特征值都等于 $1$ 时可能, 这进而蕴含 $x = 1$, 因为 $x$ 可对角化.
由习题 \ref{ex:I.5.5.4}(f), 我们知道 $\exp(\Omega)$ 包含 $1$ 在 $\UU(n)$ 中的开邻域 $V$. 由于 $\UU(n)$ 紧, 连续函数 $f$ 在 $\UU(n) - V$ 上达到上确界, 且由上段它 $< 1 + 2n$. 所以 $\sup_{x \in \UU(n)-\exp(\Omega)}|f(x)| < 1 + 2n$, 这蕴含所要的结果.
(d). 对每个 $x \in G$, 有 $|g(x)| \leq \dim W$. 所以可以取 $C = \operatorname{vol}(G - U)(\dim W)$ 并应用 (c).
(e). 设 $i\theta_1, \ldots, i\theta_n$ 是 $X$ 的特征值. 由于 $X$ 与 $X^* = -X$ 交换, 存在 $A \in \UU(n)$ 使得 $AXA^{-1} = D$, 其中 $D$ 是对角线元素为 $i\theta_1, \ldots, i\theta_n$ 的对角矩阵. 则 $Ae^XA^{-1} = e^D$ 是对角线元素为 $e^{i\theta_1}, \ldots, e^{i\theta_n}$ 的对角矩阵, 所以 $e^X$ 的特征值是 $e^{i\theta_1}, \ldots, e^{i\theta_n}$, 且
\[
f(e^X) = 1 + 2\sum_{j=1}^n\cos(\theta_j) = 1 + 2n - \sum_{j=1}^n\theta_j^2 + O\Bigl(\sum_{j=1}^n\theta_j^4\Bigr).
\]
有 $X = A^{-1}DA$, 所以 $X^* = -X = -A^{-1}DA$, 从而 $X^*X = -A^{-1}D^2A$, 最后 $\operatorname{Tr}(X^*X) = -\operatorname{Tr}(D^2) = \sum_{j=1}^n\theta_j^2$. 所以 $f(e^X) = 1 + 2n - \sum_{j=1}^n\theta_j^2 + O(\|X\|^4)$. 另一方面,
\[
(2n+1)e^{-K(X)+O(\|X\|^4)} = (2n+1)\Bigl(1 - \frac{1}{2n+1}\|X\|^2 + O(\|X\|^4)\Bigr) = 2n+1 - \sum_{j=1}^n\theta_j^2 + O(\|X\|^4).
\]
(f). 我们先证明提示. 由习题 \ref{ex:I.5.5.4}(g), 存在 $\R$-线性映射 $u \colon L \to \End(W)$ 使得对每个 $X \in L$, $\pi(e^X) = e^{tu(X)}$. 由于 $u$ 是 $\R$-线性的, 它是 $\mathcal{C}^{\infty}$ 的, 所以映射 $U \colon L \to \CC$, $X \longmapsto g(e^X) = \operatorname{Tr}(e^{u(X)})$ 也是 $\mathcal{C}^{\infty}$ 的. 还有 $U(0) = \operatorname{Tr}(\mathrm{id}_W) = \dim W$. 所以得到 $U(X) = \dim W + O(\|X\|)$.
现在我们计算积分. 做代换 $Y = N^{1/2}X$ (并注意 $NK(X) = K(Y)$), 得到
\[
\int_{\Omega}e^{-NK(X)}U(X)\,\dd X = N^{-\frac{1}{2}\dim L}\int_{N^{1/2}\Omega}e^{-K(Y)}U(N^{-1/2}Y)\,\dd Y.
\]
还需证明 $N^{1/2}\Omega$ 上的积分 $\int e^{-K(Y)}U(N^{-1/2}Y)\,\dd Y$ 在 $N \to +\infty$ 时收敛到一个正实数. 首先注意, 由于 $\int_Le^{-K(Y)}\,\dd Y$ 收敛且函数 $g$ 被 $\dim W$ 控制, 有
\[
\Bigl|\int_{N^{1/2}\Omega - N^{1/4}\Omega}e^{-K(Y)}U(N^{-1/2}Y)\,\dd Y\Bigr| \leq \int_{L - N^{1/4}\Omega}e^{-K(Y)}U(N^{-1/2}Y)\,\dd Y \xrightarrow[N \to +\infty]{} 0.
\]
另一方面, 用 $Y \in N^{1/4}\Omega$ 时 $U(N^{-1/2}Y) = \dim W + O(N^{-1/4})$ 这一事实, 得到
\begin{align*}
\lim_{N \to +\infty}\int_{N^{1/4}\Omega}e^{-K(Y)}U(N^{-1/2}Y)\,\dd Y &= \lim_{N \to +\infty}(\dim W)\int_{N^{1/4}\Omega}e^{-K(Y)}\,\dd Y \\
&= (\dim W)\int_Le^{-K(Y)}\,\dd Y.
\end{align*}
由于函数 $Y \longmapsto e^{-K(Y)}$ 在 $L$ 上取正值, 最后一个积分正且实.
(g). 若存在 $L$ 中 $0$ 的邻域 $\Omega$ 使得 $\exp$ 是从 $\Omega$ 到 $U$ 的微分同胚 (通过把 $U$ 取得足够小总可以假设), 则
\[
\int_U f(x)^Ng(x)\,\dd x = (2n + 1)^N\int_{\Omega}e^{-NK(X) + NO(\|X\|^4)}g(e^X)h(X)\,\dd X,
\]
其中 $h(X) = c + O(\|X\|)$, $c \in \R_{>0}$. 这等于
\[
\frac{(2n + 1)^N}{N^{\dim L/2}}\int_{N^{1/2}\Omega}e^{-K(Y) + O(N^{-1}\|Y\|^4)}U(N^{-1/2}Y)h(N^{-1/2}Y)\,\dd Y.
\]
我们可以像 (f) 中那样证明: 若选择 $\Omega$ 是以 $0$ 为中心的球 (我们可以), 则当 $N \to +\infty$ 时积分收敛到 $c(\dim W)\int_Le^{-K(Y)}\,\dd Y$, 这给出结论.
(h). 由 (d) 和 (f), 可以把 $\int_G f(x)^Ng(x)\,\dd x$ 分解成两项的和, 其中一项等价于 $\frac{(2n+1)^N}{(\dim L)N/2}$ 的正倍数, 另一项被 $(1 + 2n - \delta)^N$ 控制 ($\delta > 0$). 当 $N$ 趋于 $+\infty$ 时, 第二项相对于第一项变得可忽略, 所以对足够大的 $N$, 两项的和不能为 $0$.
\end{solution}

\begin{exercise}[矩阵系数方法的推广]
\label{ex:IV.9.4}
本题的目标也是把习题 \ref{ex:IV.9.2} 推广到具有忠实有限维表示 $(\rho, V)$ 的紧群 $G$. 像第 IV.9.3 节开头那样, 可以假设 $G$ 是 $\UU(N)$ 的闭子群, $V = \CC^N$, $\rho \colon G \to \UU(N) \subset \GL_N(\CC)$ 是包含.
回忆 $G$ 的所有矩阵系数的代数 $E \subset \mathcal{C}(G)$ 在第 IV.4.1 节中定义. 像第 IV.4 节那样, 把 $\mathcal{C}(G)$ 和 $E$ 看作 $G \times G$ 的表示, 让 $(x, y) \in G \times G$ 按 $L_x \circ R_y$ 作用.
对 $i, j \in \{1, \ldots, N\}$, 记 $e_{ij} \colon G \to \CC$ 为把 $x \in G \subset \GL_N(\CC)$ 送到其 $(i, j)$ 项的函数.
\begin{enumerate}
\item[(a)] 设 $E'$ 是由 $1$ 和函数 $e_{ij}$ 与 $\overline{e_{ij}}$ ($1 \leq i, j \leq N$) 生成的 $E$ 的子代数. 证明 $E'$ 被 $G \times G$ 的作用稳定.
\item[(b)] 证明存在 $\widehat{G}$ 的子集 $A$ 使得 $E' = \bigoplus_{\pi \in A}E_{\pi}$.
\item[(c)] 证明 $E' = E$.
\item[(d)] 证明 $G$ 的每个不可约表示都是形如 $V^{\otimes n} \otimes (V^*)^{\otimes m}$ (对某个 $n, m \geq 1$) 的表示的直和项.
\end{enumerate}
\end{exercise}

\begin{solution}
(a). 设 $x \in G$, 写 $x^{-1} = (a_{ij})_{1 \leq i,j \leq N}$. 则对所有 $i, j \in \{1, \ldots, N\}$, $L_xe_{ij} = \sum_{k=1}^N a_{ik}e_{kj}$, 所以 $L_xe_{ij} \in E'$. $E'$ 被算子 $R_y$ 稳定的证明类似.
(b). 回忆由定理 \ref{thm:IV.3.8}, $E = \bigoplus_{\pi \in \widehat{G}}E_{\pi}$, 且若写 $\pi \colon G \to \GL(V_{\pi})$, 则 $E_{\pi}$ 作为 $G \times G$ 的表示同构于 $\End(V_{\pi})$. 我们证明 $E_{\pi}$ 作为 $G \times G$ 的表示不可约且两两不同构. 设 $\pi, \pi' \in \widehat{G}$. 则对 $(x, y) \in G \times G$, 由命题 \ref{prop:IV.5.4} 有 $\chi_{E_{\pi}}(x, y) = \chi_{\pi}(y)\chi_{\pi}(x)$, 对 $\pi'$ 类似, 所以
\[
\langle \chi_{E_{\pi}}, \chi_{E_{\pi'}} \rangle_{L^2(G \times G)} = \int_{G^4}\chi_{\pi}(x_1)\chi_{\pi}(x_2)\overline{\chi_{\pi'}(x_3)}\overline{\chi_{\pi'}(x_4)}\,\dd x_1\,\dd x_2\,\dd x_3\,\dd x_4.
\]
由特征标的 Schur 正交性 (推论 \ref{cor:IV.5.8}(ii)), 若 $\pi \not\simeq \pi'$ 这是 $0$; 由推论 \ref{cor:IV.5.8}(iii), 若 $\pi \simeq \pi'$ 这是 $1$. $E_{\pi}$ 的不可约性由此和推论 \ref{cor:IV.5.8}(iii) 得出, $\pi \not\simeq \pi'$ 时 $E_{\pi} \not\simeq E_{\pi'}$ 也由此得出.
对每个 $\pi \in \widehat{G}$, 设 $W_{\pi}$ 是 $E'$ 在正交投影 $L^2(G) \to E_{\pi}$ 下的像 $W$. 由于 $E'$ 被 $G \times G$ 的作用稳定, $W_{\pi}$ 也如此, 所以 $W$ 是 $\{0\}$ 或 $E_{\pi}$, 因为 $E_{\pi}$ 是 $G \times G$ 的不可约表示. 所以, 若 $A = \{\pi \in \widehat{G} \mid W_{\pi} \neq \{0\}\}$, 则 $E' = \bigoplus_{\pi \in A}E_{\pi}$.
(c). $E$ 的子代数 $E'$ 被复共轭稳定, 显然分离 $G$ 上的点, 且包含常值函数 $1$, 所以由 Stone-Weierstrass 定理它在 $\mathcal{C}(G)$ 中稠密. 特别地, 它在 $E$ 中稠密. 但若 $A \subset \widehat{G}$ 是 (b) 的子集, 有 $E' = \bigoplus_{\pi \in A}E_{\pi}$, 所以 $E$ 在 $E$ 中稠密这一事实蕴含 $A = \widehat{G}$, 即 $E' = E$.
(d). 设 $W$ 是 $G$ 的不可约表示, 假设它不是任何形如 $V^{\otimes n} \otimes (V^*)^{\otimes m}$ ($n, m \geq 1$) 的表示的直和项. 由习题 \ref{ex:IV.9.2}(a), 这蕴含 $W$ 不是任何 $(\mathbf{1} \oplus V \oplus V^*)^{\otimes n}$ ($n \geq 1$) 的直和项, 即由 Schur 正交性 (定理 \ref{thm:IV.3.8}(i)), $\chi_W$ 与 $(\mathbf{1} \oplus V \oplus V^*)^{\otimes n}$ ($n \geq 1$) 的任何矩阵系数正交. 但 $(\mathbf{1} \oplus V \oplus V^*)^{\otimes n}$ 的矩阵系数的空间正好是 $E$ 中由至多 $n$ 个函数 $e_{ij}$ 和 $\overline{e_{ij}}$ ($1 \leq i, j \leq N$) 的乘积生成的子空间, 所以所有这些表示的矩阵系数生成的空间是 $E'$, 我们刚刚看到它等于 $E$. $\chi_W$ 不可能与 $E'$ 的每个元素正交, 所以我们得到结果.
\end{solution}
